\documentclass[11pt,a4paper]{article}
\usepackage{a4wide}
\setlength{\textheight}{23cm}
\setlength{\textwidth}{16cm}
%\usepackage[T1]{fontenc}
\usepackage[utf8]{inputenc}
\usepackage{latexsym}
\usepackage{amsmath}
\usepackage{amsthm}
\usepackage{amssymb,enumerate}
\usepackage[usenames]{color}
\usepackage{graphicx}
\usepackage{fancybox}
\usepackage[normalem]{ulem}
%\usepackage{showkeys}
\usepackage{hyperref}
%%%%%%%%%%%%%%%%%%%%%%%%%%%%%%%%%%%%%%%%%%%%%%%%%%%% Paquetes TKIZ
%\usepackage{graphicx} % Required for inserting images
\usepackage{pgfplots}
\pgfplotsset{compat=1.15}
%VOLVER A DESCOMENTAR %\usepackage{mathrsfs}
%\usetikzlibrary{arrows}
%\definecolor{qqqqff}{rgb}{0.,0.,1.}
%\definecolor{ududff}{rgb}{0.30196078431372547,0.30196078431372547,1.}
%\definecolor{ffqqqq}{rgb}{1.,0.,0.}
%\definecolor{xdxdff}{rgb}{0.49019607843137253,0.49019607843137253,1.}
%\definecolor{verdeosc}{rgb}{0,0.5,0}
%%%%%%%%%%%%%%%%%%%%%%%%%%%%%%%%%%%%%%%%%%%%%%%%%%%%%%%%%%%%%
\theoremstyle{plain}
%\newtheorem{theorem}{Theorem}[section]
%\newtheorem{lemma}[theo]{Lemma}
%\newtheorem{prop}[theo]{Proposition}
%\newtheorem{coro}[theo]{Corollary}
%\newtheorem*{coro*}{Corollary}
\newtheorem{proposition}{Proposition}[section]
\newtheorem{theorem}[proposition]{Theorem}
\newtheorem{lemma}[proposition]{Lemma}
\newtheorem{corollary}[proposition]{Corollary}


\theoremstyle{definition}
\newtheorem{definition}[proposition]{Definition}
\newtheorem{example}[proposition]{Example}
\newtheorem{examples}[proposition]{Examples}
\newtheorem{remark}[proposition]{Remark}
\newtheorem{remarks}[proposition]{Remarks}
%\numberwithin{equation}{section}


\usepackage{geometry}
\usepackage{subeqnarray}
\usepackage{enumerate}
\usepackage{color}

%%%%%%%%%%%%%%%%%  NACHO   %%%%%%%%%%
%Web:  https://tex.stackexchange.com/questions/249352/is-there-a-widecheck-like-widehat
%      https://tex.stackexchange.com/questions/100574/really-wide-hat-symbol

%\usepackage{scalerel,stackengine}
%\stackMath
%
%\newcommand\widecheck[1]{%
	%	\savestack{\tmpbox}{\stretchto{%
			%			\scaleto{%
				%				 \scalerel*[\widthof{\ensuremath{#1}}]{\kern-.7pt\bigwedge\kern-.4pt}%   %Original \scalerel*[\widthof{\ensuremath{#1}}]{\kern-.6pt\bigwedge\kern-.6pt}
				%				 {\rule[-\textheight/2]{1ex}{\textheight}}%WIDTH-LIMITED BIG WEDGE
				%			}{\textheight}%
			%		}{0.5ex}}%
	%	\stackon[1pt]{#1}{\scalebox{-1}{\tmpbox}}%
	%}
%\parskip 1ex


%%%%%%%%%%%%%%%%%%%%%%% Widecheck (Old! It is from the OFF Paper)
\makeatletter
\DeclareRobustCommand\widecheck[1]{{\mathpalette\@widecheck{#1}}}
\def\@widecheck#1#2{%
	\setbox\z@\hbox{\m@th$#1#2$}%
	\setbox\tw@\hbox{\m@th$#1%
		\widehat{%
			\vrule\@width\z@\@height\ht\z@
			\vrule\@height\z@\@width\wd\z@}$}%
	\dp\tw@-\ht\z@
	\@tempdima\ht\z@ \advance\@tempdima2\ht\tw@ \divide\@tempdima\thr@@
	\setbox\tw@\hbox{%
		 \raise\@tempdima\hbox{\scalebox{1}[-1]{\lower\@tempdima\box
				\tw@}}}%
	{\ooalign{\box\tw@ \cr \box\z@}}}
\makeatother
\renewcommand{\check}[1]{\widecheck{#1}}
\renewcommand{\widecheck}[1]{\widecheck{#1}}
%\usepackage{fancybox}
\usepackage{bbm}
%\usepackage{a4wide}
%\setlength{\textheight}{23cm}
%\setlength{\textwidth}{16cm}





\pagestyle{myheadings}
\newenvironment{proof1}{\medskip\par\noindent{\bf Proof.}}{\hfill $\Box$
\medskip\par}

%\newcommand{\RR}{\mathbb{R}}
%\newcommand{\CC}{\mathbb{C}}
%\newcommand{\NN}{\mathbb{N}}
%\newcommand{\N}{\mathbb{N}}
%\newcommand{\R}{\mathbb{R}}
%\newcommand{\C}{\mathbb{C}}
%\newcommand{\om}{\omega}
%\newcommand{\ka}{\kappa}
%\newcommand{\al}{\alpha}
%\newcommand{\be}{\beta}
%\newcommand{\Ga}{\Gamma}
%\newcommand{\bM}{\mathbb{M}}
%\newcommand{\M}{\mathbb{M}}
%\newcommand{\hM}{\widehat{\M}}
%\newcommand{\bL}{\mathbb{L}}
%\newcommand{\m}{\boldsymbol{m}}
%\newcommand{\bl}{\boldsymbol{\ell}}
%\newcommand{\mo}{\mathfrak{m}}
%\newcommand{\flo}[1]{\left\lfloor {#1} \right\rfloor}


%\newcommand{\ta}[1]{{\color{red}\sout{#1}}}
%\newcommand{\taf}[1]{{\color{red}#1}}
\newcommand{\RR}{\mathbb{R}}
\newcommand{\CC}{\mathbb{C}}
\newcommand{\NN}{\mathbb{N}}
\newcommand{\ZZ}{\mathbb{Z}}
\newcommand{\QQ}{\mathbb{Q}}
\newcommand{\DD}{\mathbb{D}}
\newcommand{\pr}{\operatorname{pr}}
\newcommand{\incl}{\operatorname{incl}}
\newcommand{\supp}{\operatorname{supp}}
\newcommand{\id}{\operatorname{id}}
\newcommand{\sym}{\operatorname{sym}}
\newcommand{\im}{\operatorname{im}}
\newcommand{\ev}{\operatorname{ev}}
\newcommand{\ins}{\operatorname{ins}}
\newcommand{\comp}{\operatorname{comp}}
\newcommand{\Diff}{\operatorname{Diff}}
\newcommand{\inv}{\operatorname{inv}}
\newcommand{\evol}{\operatorname{evol}}
\newcommand{\Fl}{\operatorname{Fl}}
\newcommand{\blue}[1]{{\color{blue}#1}}
%%%%%%%%%%%%            REPETIDO        %%%%%%%%%%%%%%%%%%%%%%%%%
\newcommand{\new}[1]{{\color{blue}#1}}
%%%%%%%%%%%%%%%%%%%%%%%%%%%%%%%%%%%%%%%%
\newcommand{\red}[1]{{\color{red}#1}}
\newcommand{\M}{\mathbf{M}}
\newcommand{\m}{\boldsymbol{m}}
\newcommand{\LL}{\mathbf{L}}
\newcommand{\W}{\mathbf{W}}
\newcommand{\G}{\mathbf{G}}
\renewcommand{\Re}{\text{Re}}
\let\on=\operatorname

\newcommand{\rb}[1]{\raisebox{1.5ex}[-1.5ex]{#1}}%Zentrieren von Tabellenzeilen!

\newenvironment{demo}[1]{\par\smallskip\noindent{\bf #1.}}{\par\smallskip}

\makeatletter
\@namedef{subjclassname@2020}{%
	\textup{2020} Mathematics Subject Classification}
\makeatother
%%%%%%%%%%%%%%%%%%%%%%%%%%%%%%%%%%%%%%%%%

\def\a{\alpha}
\def\b{\beta}
\def\o{\omega}
\def\ro{\rho}
\def\ga{\gamma}
%\def\Ga{\Gamma}
\def\eps{\varepsilon}
\def\bmu{\boldsymbol{\mu}}
\def\ba{\boldsymbol{\alpha}}
%\def\bM{\mathbb{M}}
%\def\M{\mathbb{M}}
\def\L{\mathbb{L}}
\def\bm{\boldsymbol{m}}
%\def\m{\boldsymbol{m}}
\def\Ms{{\mathbb{M}}_{q,\sigma}}

\def\en{\subseteq}
\def\parn{\par\noindent}
\newcommand{\pprec}{\prec\mathrel{\mkern-5mu}\prec}


\definecolor{azulosc}{rgb}{0.2,0.1,0.7}
%\definecolor{naranjauva}{RGB}{237,110,0}%col
%\definecolor{verdecla}{rgb}{0.1,0.8,0.2} %col3
%\definecolor{verdetlp}{RGB}{33,120,68}
\definecolor{granate}{rgb}{0.6,0,0.3} %col4
%\definecolor{azulclarouva}{RGB}{130,115,186}
%\definecolor{moradouva}{RGB}{186,31,181}
%\definecolor{moradopod}{RGB}{97,45,98} %moradopod
%\definecolor{moradomuyclaropod}{RGB}{160,129,161} %moradopod
%\definecolor{moradoclaropod}{RGB}{149,78,153} %moradopod
%\definecolor{verdeclaropod}{RGB}{151,194,184} %moradopod

%\definecolor{verdeuva}{RGB}{122,154,1}
%\definecolor{morado}{rgb}{0.7,0.2,0.7}        %col5
\definecolor{verdeosc}{RGB}{46, 139, 87} %colororiginal
%amarillo internet 153 130 0
%\definecolor{similar1}{RGB}{88,6,124}
%\definecolor{similar2}{RGB}{186,14,0}
%\definecolor{opuesto1}{RGB}{0,142,25}
%\definecolor{opuesto2}{RGB}{178,185,0}
%\definecolor{optotal}{RGB}{109,172,0}
%6DAC00
%\definecolor{amarillouva}{RGB}{172,104,43}     %col6
%\definecolor{rosapalo}{rgb}{0.83,0.48,0.38}   %col7
%\definecolor{violeta}{rgb}{0.7,0.5,0.8}       %col8
\definecolor{rojo}{RGB}{219,0,0}                %col9

% Widecheck
\makeatletter
\DeclareRobustCommand\widecheck[1]{{\mathpalette\@widecheck{#1}}}
\def\@widecheck#1#2{%
    \setbox\z@\hbox{\m@th$#1#2$}%
    \setbox\tw@\hbox{\m@th$#1%
       \widehat{%
          \vrule\@width\z@\@height\ht\z@
          \vrule\@height\z@\@width\wd\z@}$}%
    \dp\tw@-\ht\z@
    \@tempdima\ht\z@ \advance\@tempdima2\ht\tw@ \divide\@tempdima\thr@@
    \setbox\tw@\hbox{%
       \raise\@tempdima\hbox{\scalebox{1}[-1]{\lower\@tempdima\box
\tw@}}}%
    {\ooalign{\box\tw@ \cr \box\z@}}}
\makeatother
%%%%%%%%%%%%%%%%%%%%%%%%%%%%%%%%%%%%%%%%%

\begin{document}

\title{Stability under product and composition for uniform Carleman asymptotic expansions}
\author{Javier Jiménez-Garrido \and Ignacio Miguel-Cantero \and Javier Sanz \and Gerhard Schindl}
%\date{\today}

\maketitle

\begin{abstract}
We study the stability under point-wise product and under composition in Carleman classes of holomorphic functions, defined on sectors of the Riemann surface of the logarithm, and admitting a uniform asymptotic expansion with remainders controlled by a given sequence of positive real numbers $\M$. On the one hand, the well-known conditions of algebrability and Fa\`a di Bruno, imposed on the sequence $\M$, ensure the desired stability with respect to each operation in both the Roumieu and the Beurling settings. On the other hand, these conditions turn out to be necessary for the corresponding stability in the Roumieu case as long as the existence of suitable characteristic functions, in a precise sense, is guaranteed within the class. The construction of such functions rests on classical results of B. Rodríguez-Salinas, and is given in detail. Our results are inspired by, and thoroughly generalize, several partial statements by G.~Auberson and G.~Mennessier for Gevrey classes of order 1. 
\end{abstract}

\par\medskip

\noindent Key words: Uniform asymptotic expansion; algebras of functions; stability properties.
\par
\medskip
\noindent 2020 Mathematics Subject Classification Codes: Primary 30E15, 30H50; secondary  46J15.



\section{Introduction}



We denote by $\mathcal{\widetilde{A}}^{u}_{\{\M\}}(S)$ the ultraholomorphic Roumieu-type classes of functions with uniform asymptotic expansion, in a sector $S$ of the Riemann surface of the logarithm, whose remainders are bounded in terms of a weight sequence $\M = (M_j)_{j\in\mathbb{N}_0} \in \RR_{>0}^{\NN_0}$. This paper focuses on characterizing the stability of these classes under two fundamental operations: point-wise multiplication and composition. Our primary objective is to demonstrate that the stability properties of these classes can be characterized in terms of growth conditions on the defining sequence $\M$. For the problem of stability under point-wise multiplication, we show that the condition \hyperlink{s-alg}{$(\on{alg})$}, namely
$$\exists\;C\ge 1\;\forall\;j,k\in\NN_0:\;\;\;M_j M_k\le C^{j+k}M_{j+k},
$$
is sufficient to ensure that the class constitutes an algebra. Similarly, for the more complex problem of stability under composition, the \hyperlink{s-FdB}{$(\on{FdB})$} condition,
$$\exists\;C\ge 1\;\exists\;h\ge 1\;\forall\;k\in\NN \;\forall\;j_1,\dots,j_\ell\in\NN, \sum_{i=1}^{\ell}j_i=k:\;\;\;M_\ell\cdot M_{j_1}\cdots M_{j_{\ell}}\le Ch^kM_k,$$
serves as a sufficient requirement for closure. We mention that both statements apply to the Beurling-type classes as well.

The most challenging aspect of this work lies in proving that these sufficient conditions are, in fact, necessary in the Roumieu case. To achieve this characterization, we must ensure the existence of characteristic functions within the class. Following the ideas of B. Rodríguez-Salinas \cite{salinas62}, we define a function $f \in \widetilde{\mathcal{A}}^u_{\{\LL\}} (S)$ to be characteristic if its membership in a smaller class $\widetilde{\mathcal{A}}^u_{\{\M\}} (S) \subseteq \widetilde{\mathcal{A}}^u_{\{\LL\}} (S)$ implies the equality of the two spaces. As Theorem~\ref{optthm1} shows, this property is a consequence of the fact that the coefficients $(|c^f_j|)_{j\in\NN_0}$ of its asymptotic expansion are equivalent to the defining sequence~$\LL$.

The construction of these functions is achieved through the characteristic transform $\mathcal{T}_{\M}$, which allows us to modify some basic functions while maintaining precise control over the corresponding coefficients. By applying $\mathcal{T}_{\M}$ to these basic functions, we produce elements in $\widetilde{\mathcal{A}}^u_{\{\M\}} (S_{\alpha})$ whose coefficients $(|c_j^f|)_{j\in\NN_0}$ are equivalent to $\M$, where $S_{\alpha}$ stands for the unbounded sector bisected by the positive real line with opening $\pi\a$ in the Riemann surface of the logarithm. Furthermore, these coefficients are real, and we maintain precise control over their signs, which will be crucial in our reasoning. It should be emphasized that the requirements for the existence of characteristic functions are not uniform, varying significantly between sectors of small and large opening.

It is important to note that the proofs of both sufficiency and necessity rely on the combinatorial formulas for asymptotic expansions established by G. Auberson and G. Mennessier \cite{AubersonMennessier81}, which allow for the precise tracking of remainder terms and coefficient dependencies.

The stability properties of these classes have been a subject of extensive study in several related functional settings. In the ultradifferentiable context, A. Rainer and the fourth author \cite{compositionpaper, characterizationstabilitypaper} provided a complete characterization of these properties. In those works, the condition of derivation closedness, denoted \hyperlink{dc}{$(\on{dc})$}, is generally required for stability under composition; later, J. Pöschel \cite{Poschel} obtained results, in the weight sequences case, under the assumption of log-convexity \hyperlink{lc}{$(\on{lc})$} without requiring \hyperlink{dc}{$(\on{dc})$} (although some technical issue is to be fixed in his arguments regarding the composition of characteristic functions). It is worth noting that in the ultradifferentiable setting, the construction of characteristic functions is notably simpler than in the ultraholomorphic case. 

In the ultraholomorphic setting, earlier results often focused on specific weight sequences rather than general characterizations. For instance, G. Auberson and G. Mennessier \cite{AubersonMennessier81} proved the stability of Gevrey classes of order $1$ (corresponding to $M_j = j!$). Their result is not a characterization in the sense presented here, as Gevrey sequences inherently satisfy both \hyperlink{s-alg}{$(\on{alg})$} and \hyperlink{s-FdB}{$(\on{FdB})$} conditions. Furthermore, their technical approach to composition relied on specific combinatorial properties of factorials rather than the more general \hyperlink{s-FdB}{$(\on{FdB})$} condition. For the multivariate case of composition of asymptotic expansions, we refer to the work of J. Mozo \cite{MozoPhd}, whose results ensure stability under composition for the Gevrey classes of positive order $s$ (corresponding to $M_j = j!^s$).

Finally, in a previous work by the authors \cite{ultraholomstab}, characterization results were established for ultraholomorphic classes defined by uniform bounds on the derivatives. While those classes are closely related to the ones considered here, the arguments are fundamentally different, requiring now the construction of distinct characteristic functions and specialized technical machinery to handle the uniform nature of the asymptotic expansions.

It should be highlighted that the aforementioned results by A. Rainer and the fourth author were established for ultradifferentiable classes defined in terms of a weight matrix. This approach provides a unified framework that generalizes both the weight sequence and the weight function settings. It is expected that the characterization results obtained in this paper can be extended to ultraholomorphic classes defined in terms of weight matrices, considered in \cite{sectorialextensions, sectorialextensions1}. This generalization is currently part of an ongoing work in progress.


The structure of this work is organized as follows. We begin with a review of weight sequences and their properties, followed by the definition of the ultraholomorphic classes and the construction of the characteristic functions mentioned above. In Section~\ref{sec.formula}, we present the auxiliary formulas for the expansion of products and compositions. These results lead to Section~\ref{sec.alg}, where we characterize the closure under point-wise multiplication via the \hyperlink{s-alg}{$(\on{alg})$} condition. Finally, we address the closure under composition, establishing the equivalence between this stability property and the \hyperlink{s-FdB}{$(\on{FdB})$} condition.









%\red{When dealing with function spaces (usually called classes) it is very interesting to decide whether the usual operations (pointwise product, composition, algebraic inversion, differentiation, integration, etc.) on the functions of the space provide new functions inside it. These stability properties play a crucial role in the setting and the solution of, for example, algebraic, differential or integro-differential equations in the class.
%
%In the literature one can frequently find the so-called ultradifferentiable classes, both in the Carleman and the Braun-Meise-Taylor sense, whose elements are smooth functions defined on open subsets of $\RR^n$ (or possibly germs at a point) such that the rate of growth of their successive derivatives is controlled (except for a geometric factor) in terms of a given sequence of positive real numbers in the first case, or of a given weight function in the second one. Moreover, depending on the choice of a universal or existential quantifier for the geometric factor in the estimates, one can consider Beurling- or Roumieu-like classes in both situations. The study of stability under inversion (or division) in these frameworks has a long history, see the works of Rudin~\cite{RudinDivision}, Bruna~\cite{BrunaInverse} and Siddiqi~\cite{Siddiqi}, and also composition has been studied in Fern\'andez and Galbis~\cite{FernandezGalbis06}. Recently, the introduction of classes associated with a weight matrix, by the fourth author of this paper~\cite{diploma,dissertation}, which strictly encompass those classes mentioned before, has led him and Rainer~\cite{compositionpaper,characterizationstabilitypaper}
%to the characterization of stability under different operations in terms of conditions for the weight matrix under consideration, so giving a satisfactory general solution to these problems.
%
%In connection with the asymptotic theory of solutions for differential and difference equations around singular points in the complex domain, it is natural to consider the complex analogue of such classes, usually called ultraholomorphic classes. They consist of holomorphic functions in sectorial regions in the Riemann surface of the logarithm (the singular point is assumed to be at 0, the vertex of the region) whose derivatives admit again suitable estimates of Roumieu type in terms of a sequence of positive real numbers, which in the applications is typically a Gevrey sequence $(p!^{a})_{p\in\NN_0}$ for some $a>1$. The study of stability properties in such classes is well-known for the Gevrey ones, see~\cite{BalserUTX}, but already in 1987 Ider and Siddiqi~\cite{IderSiddiqi} studied stability under composition with analytic functions and under inversion for general Carleman-Roumieu classes in unbounded sectors not wider than a half-plane. Our aim is to extend their results in several senses: (1) we consider Roumieu classes defined by weight matrices, so including in our considerations those of Carleman type and those defined by a weight function, as in the ultradifferentiable setting; (2) we are able to deal with classes defined in sectors of arbitrary opening in the Riemann surface of the logarithm, and (3) we extend the list of stability properties, including that of composition closedness.
%It is important to note that, in the case of classes given by a weight function, a fundamental role in the stability properties is played by the condition that this function is equivalent to a concave weight function, what amounts to the root almost increasing property for the associated weight matrix.
%
%The main novelties arise from two different sources. On the one hand, the techniques coming with the weight matrix structure allow for a better understanding of the conditions usually appearing in such stability results, and provide a clear way to establish results for the weight sequence and weight function approach. Indeed, our results extend the known ones for Carleman classes, and they match, in the limit when the opening of the sector tends to 0, with the ones for ultradifferentiable classes on a half-line. On the other hand, the main statements heavily rest on the construction of so-called characteristic functions in Carleman-Roumieu ultraholomorphic classes in sectors of arbitrary opening. These functions are those in a class which cannot belong to a class strictly contained in the original one, and so are in a sense maximal within the class. While Ider and Siddiqi only got such functions in narrow sectors, the work of Rodríguez Salinas~\cite{salinas62} provides indeed the key facts for working in general sectors, and this is in turn crucial for our purposes.
%
%The paper is organized as follows. Section~\ref{sect.WeightCondit} contains all the preliminary information about sequences, weight functions and weight matrices. For the ultraholomorphic classes introduced in Section~\ref{sect.UltraholomClass} we show how to construct characteristic functions in Section~\ref{sect.CharactFunctions}. The stability results for classes associated with weight matrices are given in Section~\ref{sect.StabilityPropert}, and Section~\ref{sect.StabilityWeightFUnctCase} is devoted to their particularization to the case of classes induced by a weight function.
%{Finally, we present in Section~\ref{sect.Examples} some examples, including those of Gevrey and $q$-Gevrey classes, in order to illustrate the obtained results.}}

\section{Weight sequences}\label{sect.WeightCondit}
%\subsection{General notation}
 %With $[\cdot]$ we denote either the Roumieu case $\{\cdot\}$ or the Beurling case $(\cdot)$ but not mixing the cases (for function spaces, classes of weighted sequences and conditions).

%\subsection{Weight sequences}\label{ssectweightsequences}

We write $\NN_0:=\{0,1,2,\dots\}$ and $\NN:=\{1,2,3,\dots\}$. In what follows, we always denote by $\M=(M_j)_j\in\RR_{>0}^{\NN_0}$ a sequence with $M_0=1$.
We also use 
%$\check{\M}=(\check{M}_j)_j$ defined by $\check{M}_j:=\frac{M_j}{j!}$ and 
the sequence of quotients $\m=(m_j)_j$ defined by $m_j:=M_{j+1}/M_{j}$, $j\in \NN_0$.
%, and analogously for all other arising sequences. 
%$\M$ is called {\itshape normalized} if $1=M_0\le M_1$ holds true.\vspace{6pt}

$\M$ is said to be {\itshape log-convex}, (for short, \hypertarget{lc}{$(\text{lc})$}) if
$$\forall\;j\in\NN:\;M_j^2\le M_{j-1} M_{j+1},$$
equivalently if $\m$ is nondecreasing. If $\M$ is log-convex, %and normalized, 
then 
%both $j\mapsto M_j$ and 
$j\mapsto(M_j)^{1/j}$ is nondecreasing and $(M_j)^{1/j}\le m_{j-1}$ for all $j\in\NN$. Finally,
\begin{equation}\label{eq.alg.from.lc}
M_j M_k\le M_{j+k},\quad j,k\in\NN_0.
\end{equation}

% If $\check{\M}$ is log-convex, then $\M$ is called {\itshape strongly log-convex}, denoted by \hypertarget{slc}{$(\text{slc})$}. 
We say that a sequence $\M$ is a \emph{weight sequence} if it is \hyperlink{lc}{$(\text{lc})$} and $\lim_{j\to\infty} m_j =\infty$.
% We see that $\M$ is a \emph{normalized} weight sequence if and only if $1\leq m_0\le m_1\le\dots$, $\lim_{j\rightarrow+\infty}m_j=+\infty$ (e.g. see \cite[p. 104]{compositionpaper}) and there is a one-to-one correspondence between $\M$ and $\m$ by taking $M_j:=\prod_{i=0}^{j-1}m_i$.\vspace{6pt}

We shall use the following conditions on  sequences $\M$:

\begin{itemize}
	\item[\hypertarget{s-alg}{$(\on{alg})$}] $\M$ {is} {\itshape algebraic,} if
$$\exists\;C\ge 1\;\forall\;j,k\in\NN_0:\;\;\;M_j M_k\le C^{j+k}M_{j+k}.$$
	%\item[$\wlc$] $\M$ is \emph{weakly log-convex} if $\widehat{\M}$ satisfies $\lc$.
% 	\item[\hypertarget{s-salg}{$(\on{salg})$}] $\M$ is \emph{shifted algebraic} if
% $$\exists\;B\ge 1\;\forall\;j,k\ge 2:\;\;\;M_jM_k\le BM_{j+k-2}.$$
	% \item[\hypertarget{s-pri}{$(\on{pri})$}] $\M$ has the \emph{perturbed root increasing property,} if $$\exists\;H\ge 1\;\forall\;j\in\NN:\;\;\;(M_jH)^{1/j}\le(M_{j+1}H)^{1/(j+1)}.$$
%	\item[$\mg$] $\M$ has \emph{moderate growth} if $M_{p+q} \leq C_0H^{p+q}M_pM_q$, $p,q \in \N_0$, for some $C_0>0$ and $H>1$.
%	\item[$\nq$] $\M$ is \emph{non-quasianalytic} if $\displaystyle \sum_{p=0}^\infty \frac{1}{(p+1)m_p} < \infty$.
%	\item[$\snq$] $\M$ is \emph{strongly non-quasianalytic} if $\displaystyle \sum_{q=p}^\infty \frac{1}{(q+1)m_q} \leq  \frac{C}{m_p}$, $p \in \N_0$, for some $C > 0$.

\item[\hypertarget{s-FdB}{$(\on{FdB})$}] $\M$ has the {\itshape Fa\`{a}-di-Bruno property}, if
$$\exists\;C\ge 1\;\exists\;h\ge 1\;\forall\;k\in\NN \;\forall\;j_1,\dots,j_\ell\in\NN, \sum_{i=1}^{\ell}j_i=k:\;\;\;M_\ell\cdot M_{j_1}\cdots M_{j_{\ell}}\le Ch^kM_k.$$
% where $\check{\M}^{\circ}:=(\check{M}_j^{\circ})_{j\in\NN_0}$ is the sequence defined by
% \begin{equation*}\label{circsequ}\check{M}_k^{\circ}:=\max\left\{\check{M}_\ell\cdot \check{M}_{j_1}\cdots \check{M}_{j_{\ell}}: j_i\in\NN, \sum_{i=1}^{\ell}j_i=k\right\},\ k\in\NN;\;\;\;\;\check{M}_0^{\circ}:=1.
% \end{equation*}
% {\color{olive} $\M$ or $\check{\M}$ has FdB??}
% {\color{purple} Dado que la condición se expresa de forma natural en términos de la $\check{\M}$, me parece más correcto decir que $\check{\M}$ satisface la condición (y de hecho, diría todo en términos de una sucesión $\M$, sin mencionar $\check{\M}$). Por otra parte, en el caso de clases con desarrollo asintótico cuyos restos están controlados por una sucesión $\M$, la condición FdB hay que pedirla sobre la propia $\M$, así que todo queda muy natural. No casa tan bien cuando $\M$ controla las derivadas, pero yo en ese caso el resultado lo enunciaría (aunque no lo hemos hecho así antes): la clase $\mathcal{A}_{\{\M\}}(S)$ es estable por composición si $\check{\M}$ satisface FdB.
% }
% {\color{olive} Javier J: Me parece bien la propuesta.}
% {\color{teal} Nacho: A mi también.}
\end{itemize}

\begin{remark}\label{newcondrem}
Each of these two conditions holds when $\mathbf{M}$ is log-convex (and $M_0=1$). Indeed, from~\eqref{eq.alg.from.lc} we have \hyperlink{s-alg}{$(\on{alg})$} with $C=1$, and \cite[Lemma 2.2 (1)]{compositionpaper} shows how to obtain \hyperlink{s-FdB}{$(\on{FdB})$} from \hyperlink{lc}{$(\on{lc})$}.
% and \hyperlink{s-pri}{$(\on{pri})$} with $H=1$ by the comments above and \hyperlink{s-salg}{$(\on{salg})$} follows with $B=M_2$ by $(i)\Rightarrow(iv)$ in Theorem \ref{aubsersonmennessiercondition} below.
\end{remark}

For $a\in\RR$ we set
$$\G^a:=(j!^a)_{j\in\NN_0},\hspace{15pt}\overline{\G}^a:=(j^{ja})_{j\in\NN_0},$$
i.e. for $a>0$ the sequence $\G^a$ is the Gevrey-sequence of index $a$. Clearly $ \G^a$ and $\overline{\G}^a $ are weight sequences for any $a>0$ (by the convention $0^0:=1$).

%When $a=0$, then $G^0\equiv\overline{G}^0\equiv (1)_{j\in\NN_0}$.
%For any $M\in\hyperlink{LCset}{\mathcal{LC}}$ and $C\in\NN_{>0}$ we get that
%\begin{equation}\label{rootgeneralizedincr}
%\forall\;p\in\NN_0:\;\;\;(M_{Cp})^{1/C}\le(M_{(C+1)p})^{1/(C+1)},
%\end{equation}
%because this is equivalent to having $M_{Cp}(M_{Cp})^{1/C}\le M_{Cp+p}$ and so to $\mu_1\cdots\mu_{Cp}=M_{Cp}\le(\mu_{Cp+1}\cdots\mu_{Cp+p})^C$ which is satisfied since $p\mapsto\mu_p$ is nondecreasing.\vspace{6pt}




% $\M$ satisfies the condition of {\itshape moderate growth}, denoted by \hypertarget{mg}{$(\text{mg})$}, if
% $$\exists\;C\ge 1\;\forall\;j,k\in\NN_0:\;M_{j+k}\le C^{j+k} M_j M_k.$$
% In the classical work of Komatsu~\cite{Komatsu73} this condition is named $(M.2)$ and also known in the literature under the name {\itshape stability under ultradifferential operators}. 
$\M$ satisfies {\itshape derivation closedness}, denoted by \hypertarget{dc}{$(\text{dc})$}, if
$$\exists\;D\ge 1\;\forall\;j\in\NN_0:\;M_{j+1}\le D^{j+1} M_j\Longleftrightarrow m_{j}\le D^{j+1}.$$
In \cite{Komatsu73} this is condition $(M.2)'$. 
% Both \hyperlink{mg}{$(\on{mg})$} and \hyperlink{dc}{$(\on{dc})$} are preserved when multiplying or dividing $\M$ by any sequence $\G^a$. In particular, both conditions hold simultaneously true or false for $\M$ and $\check{\M}$.

%\vspace{6pt}
%
%We say $\M$ {has the} {\itshape root almost increasing {property}}, denoted by \hypertarget{rai}{$(\on{rai})$}, if
%the sequence of roots $({\check{M}}_j^{1/j})_{j\in\NN}$ is almost increasing, that is,
%$$\exists\;C>0\;\forall\;1\le j\le k:\;\;\;\check{M}_j^{1/j}\le C\check{M}_k^{1/k}.$$

%
%$M$ is called {\itshape nonquasianalytic,} denoted by \hypertarget{mnq}{$(\text{nq})$}, if $M$ satisfies
%$$\sum_{j=1}^{+\infty}\frac{1}{\mu_j}<+\infty.$$

%We say that $M$ has condition \hypertarget{beta1}{$(\beta_1)$} (introduced in \cite{petzsche}), if
%$$\exists\;Q\in\NN_{\ge 2}:\;\;\;\liminf_{j\rightarrow+\infty}\frac{\mu_{Qj}}{\mu_j}>Q.$$
%In \cite{petzsche} the surjectivity of the {\itshape Borel mapping} $f\mapsto(f^{(j)}(0))_{j\in\NN_0}$ from $\mathcal{E}_{[M]}$ onto the corresponding weighted sequence space has been characterized in terms of this condition. More precisely, there for $M\in\hyperlink{LCset}{\mathcal{LC}}$ it has been shown that \hyperlink{beta1}{$(\beta_1)$} is equivalent to requiring \hyperlink{gamma1}{$(\gamma_1)$}, i.e.
%$$\sup_{j\in\NN_{>0}}\frac{\mu_j}{j}\sum_{k\ge j}\frac{1}{\mu_k}<+\infty.$$
%In the literature \hyperlink{gamma1}{$(\gamma_1)$} is also called ''strong nonquasianalyticity condition'' and in \cite{Komatsu73} it is denoted by $(M.3)$ (in fact, there $\frac{\mu_j}{j}$ is replaced by $\frac{\mu_j}{j-1}$ for $j\ge 2$ but which is equivalent to having \hyperlink{gamma1}{$(\gamma_1)$}).\vspace{6pt}

Let $\M,\LL\in\RR_{>0}^{\NN_0}$ be given {with arbitrary $M_0,L_0>0$}, we write $\M\hypertarget{preceq}{\subset}\LL$ if 
$$\sup_{j\in\NN}\left(M_j/L_j\right)^{1/j}<+\infty
$$ 
or, equivalently, if there exist $A,B>0$ such that $M_j\le AB^jL_j$ for every $j\in\NN_0$. We say $\M$ and $\LL$ are {\itshape equivalent}, denoted by $\M\hypertarget{approx}{\approx}\LL$, if $\M\hyperlink{preceq}{\subset}\LL$ and $\LL\hyperlink{preceq}{\subset}\M$.
Note that, in case $M_0=L_0=1$, equivalence amounts to $B^jM_j\le L_j\le C^jM_j$ for every $j\in\NN_0$ and suitable $B,C>0$.
% Properties \hyperlink{mg}{$(\on{mg})$} and \hyperlink{dc}{$(\on{dc})$} are clearly preserved under \hyperlink{approx}{$\approx$}. 
%and for \hyperlink{beta1}{$(\beta_1)$} this follows by the characterizations obtained in \cite{petzsche}.

%Let us write $\M\le \LL$ if $M_j\le L_j$ for all $j\in\NN_0$.\vspace{6pt}

Finally{,} we recall some useful elementary estimates,
%(a consequence of) Stirling's formula
\begin{equation}\label{Stirling}
	\forall\;j\in\NN:\;\;\;\frac{j^j}{e^j}\le j!\le j^j,
\end{equation}
which immediately imply that  %the following:
%
%\begin{lemma}\label{overlineGequiv}
%	We have
 $\G^a\hyperlink{approx}{\approx}\overline{\G}^a$ for any $a\in\RR$.
%\end{lemma}

% \subsection{Associated weight function}\label{ssectAssocWeightFunct}
% Let $\M\in\RR_{>0}^{\NN_0}$, then the {\itshape associated function} $\omega_\M: \RR_{\ge 0}\rightarrow\RR\cup\{+\infty\}$ is defined by
% \begin{equation*}\label{assofunc}
% \omega_\M(t):=\sup_{j\in\NN_0}\ln\left(\frac{t^j}{M_j}\right)\;\;\;\text{for}\;t>0,\hspace{30pt}\omega_\M(0):=0.
% \end{equation*}
% For an abstract introduction of the associated function we refer to \cite[Chapitre I]{mandelbrojtbook}, see also \cite[Definition 3.1]{Komatsu73}. %Note that $\omega_M$ is here extended to whole $\RR$ in an symmetric (even) way.

% If $\liminf_{j\rightarrow+\infty}(M_j)^{1/j}>0$, then $\omega_\M(t)=0$ for sufficiently small $t>0$, since $t^0/M_0=1$ and $\ln\left(\frac{t^j}{M_j}\right)<0$ precisely if $t<(M_j)^{1/j}$, $j\in\NN$ (in particular, if $M_j\ge 1$ for all $j\in\NN_0$, then $\omega_\M$ vanishes on $[0,1]$). Moreover, under this assumption $t\mapsto\omega_\M(t)$ is a continuous nondecreasing function, which is convex in the variable $\ln(t)$ and tends faster to infinity than any $\ln(t^j)$, $j\ge 1$, as $t\rightarrow+\infty$. If $\lim_{j\rightarrow+\infty}(M_j)^{1/j}=+\infty$, then $\omega_\M(t)<+\infty$ for each finite $t$, so this will be a basic assumption for defining $\omega_\M$.\vspace{6pt}

% If $\M$ is a weight sequence, then we can compute $\M$ by involving $\omega_\M$ as follows, see \cite[Chapitre I, 1.4, 1.8]{mandelbrojtbook} and also \cite[Prop. 3.2]{Komatsu73}:
% \begin{equation}\label{Prop32Komatsu}
% 	M_j=\sup_{t\ge 0}\frac{t^j}{\exp(\omega_{\M}(t))},\;\;\;j\in\NN_0.
% \end{equation}
% Moreover, in this case one has
% \begin{equation*}%\label{assovanishing}
% 	\omega_\M(t)=0\hspace{20pt}\forall\;t\in[0,m_0],
% \end{equation*}
% by the known integral representation formula for $\omega_\M$, see \cite[1.8. III]{mandelbrojtbook} and also \cite[$(3.11)$]{Komatsu73}.

% If $\M\in\RR_{>0}^{\NN_0}$ satisfies $\lim_{j\rightarrow+\infty}(M_j)^{1/j}=+\infty$, then the right-hand side of formula \eqref{Prop32Komatsu} yields the $j$-th term of the log-convex minorant $\M^{\on{lc}}$ of $\M$, i.e. the log-convex sequence such that each log-convex sequence $\LL$ with $\LL\le \M$ satisfies $\LL\le \M^{\on{lc}}$ (moreover, $\M^{\on{lc}}\equiv \M$ if and only if $\M$ is log-convex). By the results from \cite[Chapitre I]{mandelbrojtbook} it also follows that $\omega_\M\equiv\omega_{\M^{\on{lc}}}$.

% Finally, if for $\beta>0$ we write $\M^{1/\beta}:=(M_j^{1/\beta})_{j\in\NN_0}$, we recall the following immediate equality, e.g. see  \cite[$(2.7)$]{sectorialextensions}:
% \begin{equation*}%\label{assofctpower}
% 	\forall\;t\ge 0:\;\;\;\omega_{\M}^{\beta}(t):=\omega_{\M}(t^{\beta})=\beta\omega_{\M^{1/\beta}}(t).
% \end{equation*}

%\subsection{Growth index $\gamma(\M)$}\label{ssect.IndexgammaM}
%
%We say $\M$ satisfies property $\left(P_{\gamma}\right)$ if there exists a sequence of real numbers $\boldsymbol{\ell}=(\ell_{j})_{j\in\NN_0}$
%such that:
%\begin{enumerate}[(i)]
%	\item $\m\simeq\boldsymbol{\ell}$, that is,
%	$$\exists\;a\ge 1\;\forall\;j\in\NN_0:\;\;\;a^{-1}m_j\le \ell_j\le a m_j,$$
%	
%	\item $\left((j+1)^{-\gamma}\ell_{j}\right)_{j\in\NN_0}$ is nondecreasing.
%\end{enumerate}
%Note that $\m\simeq\boldsymbol{\ell}$ implies $\M\hyperlink{approx}{\approx}\LL$.
%
%If $(P_\gamma)$ holds true for $\M$, then $(P_{\gamma'})$ also holds for any $\gamma'\leq\gamma$.
%It is then natural to define the \textit{growth index} $\gamma(\M)$ by
%$$\gamma(\M):=\sup\{\gamma\in\RR:\, (P_{\gamma})\hbox{ is fulfilled}\},$$%
%with the conventions $\inf\emptyset=\sup\RR=+\infty$ and $\inf\RR=\sup\emptyset=-\infty$ (see \cite[Rem. 2.2]{index}). For a comprehensive study of this index we refer to \cite[Sect. 3]{index}, especially to the characterizing result \cite[Thm. 3.11]{index}.
%This growth index was originally defined and considered for so-called {\itshape strongly regular sequences} by V. Thilliez in \cite[Sect. 1]{Thilliezdivision}.





\section{Uniform asymptotic expansion classes}\label{sect.AsymptoticexpClass}

We introduce now the crucial classes under consideration in this paper, i.e., classes of functions that admit a uniform asymptotic expansion at the vertex of the sector where they are defined. We define spaces of both Roumieu and Beurling type, analogously as it was done for classes weighting the derivatives of the functions under consideration; for the ultraholomorphic setting see~\cite{ultraholomstab}, and~\cite{compositionpaper} for ultradifferentiable classes.
% (see Section \ref{sect.UltraholomClass} for the Roumieu case and \ref{sect.Beur.UltraholomClass}; for the Beurling type).


Recall that $\mathcal{R}$ stands for the Riemann surface of the logarithm. 
We will work with sectors in $\mathcal{R}$ of the form
$$
\{z\in\mathcal{R}\colon \beta<\arg(z)<\delta,\ 0<|z|<R\},
$$
where $\beta,\delta\in\RR$ with $\beta<\delta$, and $R\in(0,\infty]$. They can be seen as part of the complex plane whenever $\delta-\beta\le 2\pi$. They are called bounded if $R\in(0,\infty)$, unbounded otherwise.

For $\a>0$, we consider unbounded sectors bisected by direction 0 and opening $\a\pi$,
$$S_{\a}:=\{z\in\mathcal{R}:|\arg(z)|<\frac{\a\,\pi}{2}\}.$$
%
Let $T$ and $S$ be sectors in $\mathcal{R}$ with vertex at $0$. We say that $T$ is a \emph{subsector} of $S$ whenever $T\subset S$, and $T$ is a \emph{proper subsector} of  $S$ if $\overline{T}\subset S$ (where the closure of $T$ is taken in $\mathcal{R}$, and so the vertex of the sector is not under consideration).

%
%or, in general, unbounded sectors with bisecting direction $d\in\R$ and opening $\ga\,\pi$,
%$$S(d,\ga):=\{z\in\mathcal{R}:|\hbox{arg}(z)-d|<\frac{\ga\,\pi}{2}\}.$$
%

%We denote by $\CC[[z]]$ the space of formal power series in $z$ with complex coefficients.
We start by recalling the concept of uniform asymptotic expansion.

Let $\M$ be a sequence, $S\subseteq\mathcal{R}$ a sector and $h>0$. We define
$\mathcal{\widetilde{A}}^{u}_{\M,h}(S)$ as the space of $f\in\mathcal{H}(S)$ for which there exists a formal complex power series $\sum_{k=0}^{\infty}c_kz^k$ such that 
\[
\left\|f\right\|_{\M,h,\overset{\sim}{u}}:=\sup_{z\in S,j\in\NN_{0}}\frac{|f(z)-\sum_{k=0}^{j-1}c_kz^k|}{h^{j}M_{j}|z|^j}<+\infty.
\]
$(\mathcal{\widetilde{A}}^{u}_{\M,h}(S),\|\cdot\|_{\M,h,\overset{\sim}{u}})$ is a Banach space and alternatively we write $f\sim^u_{\M,h} \sum_{j=0}^\infty c_j z^j$ if $f\in\mathcal{\widetilde{A}}^{u}_{\M,h}(S)$ and say that $f$ admits $\sum_{k=0}^{\infty}c_kz^k$ as its \emph{uniform $\M$-asymptotic expansion of type $1/h$.} Next we put
\begin{equation*}%\label{spacedefine}
		 \mathcal{\widetilde{A}}^{u}_{\{\M\}}(S):=\bigcup_{h>0}\mathcal{\widetilde{A}}^{u}_{\M,h}(S),
\end{equation*}
and so $\widetilde{\mathcal{A}}^u_{\{\M\}}(S)$ stands for the $(LB)$ space of functions admitting a uniform $\{\M\}$-asymptotic expansion (of Roumieu type) in $S$. Occasionally, we write $f\sim^u_{\{\M\}} \sum_{j=0}^\infty c_j z^j$ if $f\in\widetilde{\mathcal{A}}^u_{\{\M\}}(S)$. Moreover, we can consider the space of Beurling-type, denoted by $\widetilde{\mathcal{A}}^u_{(\M)}(S)$, and defined as
$$
\widetilde{\mathcal{A}}^u_{(\M)}(S):=\bigcap_{h>0}\widetilde{\mathcal{A}}^u_{\M,h}(S),
$$
which becomes a Fr\'echet space when endowed with the topology generated by the family of seminorms $(\left\|\cdot\right\|_{\M,h,\overset{\sim}{u}})_{h>0}$. In this case, we write $f\sim^u_{(\M)} \sum_{j=0}^\infty c_j z^j$ if $f\in\widetilde{\mathcal{A}}^u_{(\M)}(S)$.
%By definition it is immediate that $\M\hyperlink{approx}{\approx}\LL$ implies %$\mathcal{A}_{\{\M\}}(S)=\mathcal{A}_{\{\LL\}}(S)$ and
%$\mathcal{\widetilde{A}}^{u}_{\{\M\}}(S)=\mathcal{\widetilde{A}}^{u}_{\{\LL\}}(S)$ (as locally convex vector spaces) for any sector $S$.

Let us fix some notation. If $f\in \widetilde{\mathcal{A}}^u_{\{\M\}} (S)$  with $f\sim^u_{\{\M\}} \sum_{j=0}^\infty c_j z^j$, then we denote the remainder
\begin{equation*}%\label{remainderequ}
	r_{f}(z,n):= f(z)-\sum_{j=0}^{n-1} c_j z^j,\;\;\;z\in S,\;n\in\NN_0,
\end{equation*}
and also observe that
\begin{equation*}%\label{remainderequnew}
	\forall\;z\in S\;\forall\;n\in\NN_0:\;\;\; \frac{zr_f(z,n+1)}{z^{n+1}}=\frac{r_f(z,n)}{z^{n}}-c_n.
\end{equation*}
If $f\in\widetilde{\mathcal{A}}^u_{\{\M\}} (S)$, then $z\mapsto z^{-n}r_f(z,n)$ is bounded in $S$ for each $n\in\NN_0$ fixed. This implies
\begin{equation}\label{eq.coeff.asymp.expan.lim.remainder}
    \forall\;n\in\NN_0:\;\;\;\lim_{z\in S,z\rightarrow 0} z^{-n}r_f(z,n)=c_n.
\end{equation}
Moreover, if $f\sim^u_{\M,h} \sum_{j=0}^\infty c_j z^j$ on $S$, then by \eqref{eq.coeff.asymp.expan.lim.remainder} we have
\begin{equation}\label{cjestimation}
	\exists\;A>0\;\forall\;n\in\NN_0:\;\;\;|c_n|\le Ah^nM_n.
\end{equation}

It is also a well-known fact, steming from Cauchy's integral formula for the derivatives, that for every $n\in\NN_0$ and every proper subsector $T$ of $S$ one has
\begin{equation}\label{eq.limit.derivatives.coeff}
\lim_{\substack{z\to 0\\z\in T}}\frac{f^{(n)}(z)}{n!}=c_n.    
\end{equation}

\begin{remark}
%	First, note that, taking $p=0$ in~\eqref{desarasintunifo}, we deduce that every function in $\widetilde{\mathcal{A}}^u_{\{\M\}}(S)$ or $\widetilde{\mathcal{A}}^u_{(\M)}(S)$ is a bounded function.
%	
%	Secondly,
	When a statement is valid for both Roumieu and Beurling classes, we will use the notation $\widetilde{\mathcal{A}}^u_{[\M]}(S)$ and so on (substituting every square bracket by either of them, curly brackets or parentheses, but the same all through the statement). Moreover, if $\M\approx\LL$ the corresponding classes coincide (as locally convex vector spaces), i.e. $\widetilde{\mathcal{A}}^u_{[\M]}(S)=\widetilde{\mathcal{A}}^u_{[\LL]}(S)$.
	%if $\M$ is $\on{lc}$, the spaces $\mathcal{A}_{[\M]}(S)$ and $\widetilde{\mathcal{A}}^u_{[\M]}(S)$
	%are algebras, and if $\M$ is $\on{dc}$ they are stable under taking derivatives. Moreover,
	
\end{remark}	
%\vspace{6pt}
%
%     CLASES ASOCIADAS A FUNCIONES PESO
%

% Similarly, let $\omega\in\hyperlink{omset0}{\mathcal{W}_0}$ be a weight function, $S$ an (unbounded) sector and for every $\ell>0$, we first define
% \[\mathcal{\widetilde{A}}^{u}_{\omega,\ell}(S):=\{f\in\mathcal{H}(S): \|f\|_{\omega,\ell}:=\sup_{z\in S, j\in\NN_0}\frac{|f(z)-\sum_{k=0}^{j-1}c_kz^k|}{\exp(\frac{1}{\ell}\varphi^{*}_{\omega}(\ell j))}<+\infty\}.\]
% $(\mathcal{\widetilde{A}}^{u}_{\omega,\ell}(S),\|\cdot\|_{\omega,\ell})$ is a Banach space and we put
% \begin{equation*}%\label{spacedefine}
% 	 \mathcal{\widetilde{A}}^{u}_{\{\omega\}}(S):=\bigcup_{\ell>0}\mathcal{\widetilde{A}}^{u}_{\omega,\ell}(S),\qquad \qquad \mathcal{\widetilde{A}}^{u}_{(\omega)}(S):=\bigcap_{\ell>0}\mathcal{\widetilde{A}}^{u}_{\omega,\ell}(S).
% \end{equation*}
% $\mathcal{\widetilde{A}}^{u}_{\{\omega\}}(S)$ is the space of functions admitting a uniform $\{\omega\}$-asymptotic expansion in $S$ of Roumieu type (it is an $(LB)$ space), instead $\mathcal{\widetilde{A}}^{u}_{(\omega)}(S)$ of Beurling type (it is an $(LF)$ space). Again, equivalent weight functions provide equal associated classes: This follows by the above definition and (the proof of) \cite[Lemma 5.16 (1)]{compositionpaper}.

%
%     CLASES ASOCIADAS A MATRICES PESO
%

% Finally, we define uniform asymptotic expansion classes defined by a weight matrix $\mathcal{M}$.

% Given a weight matrix $\mathcal{M}=\{\M^{(\alpha)}\in\RR_{>0}^{\NN_0}: \alpha\in\mathcal{I}\}$ and a sector $S$ we may introduce the class $\mathcal{\widetilde{A}}^{u}_{\{\mathcal{M}\}}(S)$  of {\itshape Roumieu type} as
% \begin{equation*}%\label{generalroumieuultraholom}
% 	 \mathcal{\widetilde{A}}^{u}_{\{\mathcal{M}\}}(S):=\bigcup_{\alpha\in\mathcal{I}}\mathcal{\widetilde{A}}^{u}_{\{\M^{(\alpha)}\}}(S),
% \end{equation*}
% and the class $\mathcal{\widetilde{A}}^{u}_{(\mathcal{M})}(S)$  of {\itshape Beurling type} as
% \begin{equation*}%\label{generalroumieuultraholom}
% 	\mathcal{\widetilde{A}}^{u}_{(\mathcal{M})}(S):=\bigcap_{\alpha\in\mathcal{I}}\mathcal{\widetilde{A}}^{u}_{(\M^{(\alpha)})}(S).
% \end{equation*}
% R-equivalent (resp. B-equivalent) weight matrices yield (as locally convex vector spaces) the same function class of Roumieu type (resp. Beurling type) on each sector $S$.

% Let now $\omega\in\hyperlink{omset1}{\mathcal{W}}$ be given and let $\mathcal{M}_{\omega}$ be the associated weight matrix defined in Subsection \ref{ssect.WeightMatricesFromFunctions}, then
% \begin{equation}\label{equaEqualitySpacesWeightFunctionMatrix}
% 	 \mathcal{\widetilde{A}}^{u}_{[\omega]}(S)=\mathcal{\widetilde{A}}^{u}_{[\mathcal{M}_{\omega}]}(S)
% \end{equation}
% holds as locally convex vector spaces. This equality is an easy consequence of \eqref{newexpabsorb} and the way the seminorms are defined in these spaces. Note that we use square brackets to refer to both cases, Roumieu and Beurling. %\vspace{6pt}

% On the other hand, by $(iii)$ in Subsection \ref{ssect.WeightMatricesFromFunctions} we get the following:
% \begin{lemma}\label{om6lemma}
% 	Let $\omega\in\hyperlink{omset1}{\mathcal{W}}$ be given and assume that $\omega$ has \hyperlink{om6}{$(\omega_6)$}. Then, for all sectors $S$ we get that
% $$\forall\;\ell>0:\;\;\;\mathcal{\widetilde{A}}^{u}_{[\omega]}(S)=\mathcal{\widetilde{A}}^{u}_{[\W^{(\ell)}]}(S),$$
% as locally convex vector spaces.
% \end{lemma}


%
%    ESTO YA ESTABA COMENTADO
%
%If $f$ belongs to any of such classes, we may define the complex numbers $$f^{(j)}(0):=\lim_{z\in S, z\rightarrow 0}f^{(j)}(z),\quad j\in\NN_0.$$

%
%\section{Characteristic functions in ultraholomorphic classes}\label{sect.CharactFunctions}
%We follow the ideas from \cite[Sect. 4]{ultraholomstab} and start with the following definition; see the analogous \cite[Def. 4.1]{ultraholomstab}.
%
%\begin{definition}\label{chardef}
%	Let $\LL\in\RR_{>0}^{\NN_0}$ and $S$ be a given sector. A function $f\in \widetilde{\mathcal{A}}^u_{\{\LL\}} (S)$ is said to be \textit{characteristic} in the class $\widetilde{\mathcal{A}}^u_{\{\LL\}} (S)$ if, whenever $f\in \widetilde{\mathcal{A}}^u_{\{\M\}} (S)\subseteq \widetilde{\mathcal{A}}^u_{\{\LL\}} (S)$ for some $\M\in\RR_{>0}^{\NN_0}$, we have that $\widetilde{\mathcal{A}}^u_{\{\M\}} (S)= \widetilde{\mathcal{A}}^u_{\{\LL\}} (S)$.
%\end{definition}
%
%Let us fix some notation. If $f\in \widetilde{\mathcal{A}}^u_{\{\LL\}} (S)$  with $f\sim^u_{\{\LL\}} \sum_{j=0}^\infty c_j z^j$, then put
%\begin{equation}\label{remainderequ}
%r_{f}(z,n):= f(z)-\sum_{j=0}^{n-1} c_j z^j,\;\;\;z\in S,\;n\in\NN_0,
%\end{equation}
%and also observe that
%\begin{equation}\label{remainderequnew}
%\forall\;z\in S\;\forall\;n\in\NN_0:\;\;\;z r_f(z,n+1)/z^{n+1}=r_f(z,n)/z^{n}-c_n.
%\end{equation}
%If $f\in\widetilde{\mathcal{A}}^u_{\{\LL\}} (S)$, then $z\mapsto z^{-n}r_f(z,n)$ is bounded in $S$ for each $n\in\NN_0$ fixed. This implies
%$$\forall\;n\in\NN_0:\;\;\;\lim_{z\in S,z\rightarrow 0} z^{-n}r_f(z,n)=c_n.$$
%Moreover, if $f\sim^u_{\LL,h} \sum_{j=0}^\infty c_j z^j$ on $S$, then by \eqref{remainderequnew}.
%\begin{equation}\label{cjestimation}
%\exists\;A>0\;\forall\;z\in S\;\forall\;n\in\NN_0:\;\;\;|c_n|\le Ah^nL_n.
%\end{equation}
%Finally, we consider the sequence defined by
%\begin{equation}\label{tildesequenceC}
%	\widetilde{C}_n(f):=\sup_{z\in S}|z^{-n}r_f(z,n)|, \qquad n\in\NN_0.
%\end{equation}
%
%\begin{theorem}\label{optthm1}
%	Let $\LL\in\RR_{>0}^{\NN_0}$, $S$ be a given sector and $f\in \widetilde{\mathcal{A}}^u_{\{\LL\}} (S)$ with $f\sim^u_{\{\LL\}} \sum_{j=0}^\infty c_j z^j$. Then, each of the following  conditions implies the next one:
%	\begin{enumerate}
%		\item[$(1)$] The sequence $(|c_j|)_{j\in\NN_0}$ is equivalent to $\LL$.
%		\item[$(2)$] The sequence $(\widetilde{C}_j(f))_{j\in\NN_0}$ is equivalent to $\LL$.
%		\item[$(3)$] $f$ is characteristic in the class $\widetilde{\mathcal{A}}^u_{\{\LL\}} (S)$.
%	\end{enumerate}
%\end{theorem}
%
%\demo{Proof}
%(1) $\Rightarrow$ (2) As $f\in \widetilde{\mathcal{A}}^u_{\{\LL\}} (S)$, there exist $A,B>0$ such that $\widetilde{C}_n(f)\le AB^nL_n$ for every $n\in\NN_0$. On the other hand, it is clear that $\widetilde{C}_n(f)\geq |c_n|$, and the hypothesis allows us to conclude the other estimate.
%
%(2)  $\Rightarrow$ (3) By assumption, there exist $A,B>0$ such that $L_n\le A{B}^n\widetilde{C}_n(f)$ for every $n\in\NN_0$. If for some $\M=(M_n)_{n\in\NN_0}\in\RR_{>0}^{\NN_0}$ we have $f\in \widetilde{\mathcal{A}}^u_{\{\M\}} (S)\subseteq \widetilde{\mathcal{A}}^u_{\{\LL\}} (S)$, there exist $C,D>0$ such that $\widetilde{C}_n(f)\le C {D}^nM_n$ for every $n\in\NN_0$. The two deduced inequalities show that $L_n\le{ AC(BD)^n}M_n$ for every $n\in\NN_0$, what easily implies that $\widetilde{\mathcal{A}}^u_{\{\LL\}} (S)\subseteq \widetilde{\mathcal{A}}^u_{\{\M\}} (S)$, and we are done.
%\qed\enddemo

%\subsection{Basic functions}\label{ssect.BasicFct}
%
%Recall the notations {$\G^s:=(j!^s)_{j\in\NN_0}$ and }$\overline{\G}^s:=(j^{js})_{j\in\NN_0}$, $s\in\RR$, and that $\overline{\G}^s\hyperlink{approx}{\approx}\G^s$, see \eqref{Stirling}.\vspace{6pt}
%
%The \textit{{two-parametric} Mittag-Leffler function} is defined for all {complex parameters $A, B$ with $\Re(A)>0$} by
%$$E_{A,B}(z):=\sum_{j=0}^\infty \frac{z^j}{\Gamma (Aj+B)},\quad z\in\CC,$$
%where $\Gamma$ denotes the Gamma function.
%For the construction of characteristic functions in sectors $S_\alpha$ for $\alpha\in(0,1]$ we will take $A=2-\alpha$ and $B=4-\alpha$ and we set
%\begin{equation*}%\label{goodEalpha}
%	\widetilde{E}_{\alpha}(z):=E_{2-\alpha, 4-\alpha}(-z){=\sum_{j=0}^\infty \frac{(-1)^jz^j}{\Gamma ((2-\alpha)(j+1)+2)},\quad z\in\CC}.
%\end{equation*}
%
%We recall the following statement.
%
%\begin{theorem}\label{Ealphathm} (\cite[Thm. 5, Thm. 20]{salinas62})
%	Let $\alpha\in(0,2)$, then
%	\begin{equation}\label{Ealphagoodestimate}
%		\exists\;C\ge 1\;\forall\;z\in S_{\alpha}\;\forall\;n\in\NN_0:\;\;\;\left| \widetilde{E}_{\alpha}(z)-\sum_{j=0}^{n-1} \frac{(-z)^j}{\Gamma ((j+1) (2-\alpha)+2)}\right|\leq C \frac{e^n}{n^{(2-\alpha)n}}|z|^n.
%	\end{equation}
%	Consequently, $\widetilde{E}_{\alpha}\in\widetilde{\mathcal{A}}^u_{\{\overline{\G}^{\alpha-2}\}}(S_\alpha)$ and $\widetilde{E}_{\alpha}$ is a characteristic function in the class $\widetilde{\mathcal{A}}^u_{\{\overline{\G}^{\alpha-2}\}}(S_\alpha)$.
%\end{theorem}
%
%Next we consider the \textbf{Rodr\'{i}guez-Salinas function} (see \cite{salinas62})
%\begin{equation}\label{RSfunction}
%	K_{RS}(z):=\frac{(1+z)\log(1+z)-z}{z^2},
%\end{equation}
%which is clearly holomorphic in $S_2$.
%
%This function satisfies the following result:
%
%\begin{theorem}\label{RStheorem} (\cite[Thm. 3]{salinas62})
%	There exists a constant $C$ (we can take $C=4$) such that
%	 $$\forall\;z\in S_2
%	 \;\forall\;n\in\NN_0:\;\;\;\left| K_{RS}(z)-\sum_{j=0}^{n-1} \frac{(-z)^j}{(j+1)(j+2)}\right|\leq C|z|^n.$$
%	Consequently, for the constant sequence $\mathbbm{1}:=(1)_{j\in\NN_0}$ we have that  $K_{RS}\in\widetilde{\mathcal{A}}^u_{\{\mathbbm{1}\}}(S_2)$ and, moreover, $K_{RS}$ is a characteristic function in the class $\widetilde{\mathcal{A}}^u_{\{ \mathbbm{1}\}}(S_2)$.
%\end{theorem}
%%
%%
%Finally, let $\alpha>2$ and take $\alpha'>\alpha$. For all $z\in S_\alpha$ we define
%\begin{equation}\label{fctfalphaalphaprime}
%	f_{\alpha,\alpha'}(z):= \int_{0}^{\infty(-\phi)} K_{RS}(zv^{\alpha'-2}) e^{-v} dv,
%\end{equation}
%where we choose $\phi  \in (-\frac{(\alpha-2)}{(\alpha'-2)}\frac{\pi}{2}, \frac{(\alpha-2)}{(\alpha'-2)}\frac{\pi}{2})$ with $|\arg(z)- (\alpha'-2)\phi|<\pi$.\\
%\begin{theorem}\label{fgalphaalphaprimethm} (\cite[Thm. 4, Thm. 8]{salinas62})
%	Let $\alpha>2$, $\alpha'>\alpha$ and
%	$f_{\alpha,\alpha'}$ be the function from
%	\eqref{fctfalphaalphaprime}. Then
%				\begin{align}\label{falphagoodestimate}
%						\exists\;C,A\ge 1\;\forall\;z\in S_{\alpha}\;\forall\;n\in\NN_0:\;\;\;&\left|f_{\alpha,\alpha'} (z) -\sum_{j=0}^{n-1} \frac{(-z)^j}{(j+1)(j+2)} \Gamma (  (\alpha'-2)j+1)\right|\nonumber
%						\\&
%						\le CA^n \Gamma ( (\alpha'-2)n+1)|z|^n.
%					\end{align}
%				Consequently, $f_{\alpha,\alpha'}\in\widetilde{\mathcal{A}}^u_{\{\overline{\G}^{\alpha'-2}\}}(S_\alpha)$ and,
%				moreover, $f_{\alpha,\alpha'}$ is a characteristic function in the class $\widetilde{\mathcal{A}}^u_{\{ \overline{\G}^{\alpha'-2} \}}(S_\alpha)$.
%	\end{theorem}
%%Let $\alpha>1$ and take $\alpha'>\alpha$. For all $z\in S_\alpha$ we define
%%\begin{equation}\label{fctgalphaalphaprime}
%%	g_{\alpha,\alpha'} (z):= \int_{0}^{\infty(-\phi)} e^{-zv^{\alpha'-1}} e^{-v} dv,
%%\end{equation}
%%where we choose $\phi  \in (-\frac{(\alpha-1)}{(\alpha'-1)}\frac{\pi}{2}, \frac{(\alpha-1)}{(\alpha'-1)}\frac{\pi}{2})$ with $|\arg(z)- (\alpha'-1)\phi|<\pi/2$.
%%
%%\begin{theorem}\label{fgalphaalphaprimethm} (\cite[Thm. 28]{salinas62})
%%	Let $\alpha>1$, $\alpha'>\alpha$ and
%%%$f_{\alpha,\alpha'}$ resp.
%%$g_{\alpha,\alpha'}$ be the function from
%%%\eqref{fctfalphaalphaprime} resp.
%%\eqref{fctgalphaalphaprime}.
%%%	\begin{itemize}
%%%		\item[$(i)$] By \cite[Thm. 4, Thm. 8]{salinas62} we get: Let $\alpha>2$ and take $\alpha'>\alpha$. Then
%%%		\begin{align*}
%%%			\exists\;C,A\ge 1\;\forall\;z\in S_{\alpha}\;\forall\;n\in\NN_0:\;\;\;&\left|f_{\alpha,\alpha'} (z) -\sum_{j=0}^{n-1} \frac{(-z)^j}{(j+1)(j+2)} \Gamma (  (\alpha'-2)j+1)\right|
%%%			\\&
%%%			\le CA^n \Gamma ( (\alpha'-2)n+1)|z|^n.
%%%		\end{align*}
%%%		Consequently, we have that $f_{\alpha,\alpha'}\in\widetilde{\mathcal{A}}^u_{\{\overline{\G}^{\alpha'-2}\}}(S_\alpha)$ and,
%%%		moreover, $f_{\alpha,\alpha'}$ is a characteristic function in the class $\widetilde{\mathcal{A}}^u_{\{ \overline{\G}^{\alpha'-2} \}}(S_\alpha)$.
%%		
%%%		\item[$(ii)$]
%%Then,
%%		\begin{equation}\label{galphagoodestimate}
%%			\exists\;C,A\ge 1\;\forall\;z\in S_{\alpha}\;\forall\;n\in\NN_0:\;\;\;\left| g_{\alpha,\alpha'}^{(n)}(z)\right|\leq C A^n \Gamma ((\alpha'-1)n+1).
%%		\end{equation}
%%		Consequently, $g_{\alpha,\alpha'}\in \mathcal{A}_{\{\overline{\G}^{\alpha'-1}\}}(S_\alpha)$. Moreover,
%%{$$
%%g_{\alpha,\alpha'}^{(n)}(0)=(-1)^n  \Gamma((\alpha'-1)n+1),\qquad n\in\NN_0,
%%$$
%%and so} $g_{\alpha,\alpha'}$ is a characteristic function of the class $\mathcal{A}_{\{\overline{\G}^{\alpha'-1}\}}(S_\alpha)$.
%%%\end{itemize}
%%\end{theorem}
%
%\subsection{Characteristic transform}\label{ssect.CharactTransf}
%We revisit the notion of the recent functional transform and its consequences from \cite[Sect. 4.2]{ultraholomstab}. This transform modifies the derivatives at 0 of a given function in an ultraholomorphic class. We start with the analogue of \cite[Def. 4.5]{ultraholomstab}:
%
%\begin{definition}
%	Let $\M$ be an (lc) sequence, $\LL\in\RR_{>0}^{\NN_0}$, $S$ a sector and $f\in\widetilde{\mathcal{A}}^u_{\{\LL\}} (S)$
%	%resp. $f\in \mathcal{A}_{\{\LL\}} (S)$
%. Then we define the $\mathcal{T}_{\M}-$transform of $f$ by
%	\begin{equation*}%\label{Ttrasform}
%		\mathcal{T}_{\M}(f)(z):= \sum_{j=0}^{\infty}\frac{1}{2^j} \frac{M_j}{m_j^j} f(m_j z), \qquad z\in S.
%	\end{equation*}
%\end{definition}
%For every $j\in\NN_0$ let us set
%\begin{equation*}%\label{Sexpression}
%	R_j:=\sum^{\infty}_{n=0} \frac{1}{2^n} \frac{M_n}{m_n^n}m_n^j.
%\end{equation*}
%
%Now let us recall \cite[Lemma 4.6]{ultraholomstab}
%
%\begin{lemma}\label{Scomplemma}
%	Let $\M\in\RR_{>0}^{\NN_0}$, then
%	$$\forall\;j\in\NN_0:\;\;\;R_j\ge\frac{1}{2^j} M_j.$$
%	If $\M$ is (lc), then also
%	$$\forall\;j\in\NN_0:\;\;\;R_j\le 2 M_j,$$
%	and so $(R_j)_{j\in\NN_0}$ is equivalent to $\M$.
%\end{lemma}
%
%
%\begin{theorem}\label{CharacteristicTransformUniformAsymptotic}
%        Let $\M$ be a (lc) sequence, $\LL\in\RR_{>0}^{\NN_0}$ and for a given sector $S$ take $f\in\widetilde{\mathcal{A}}^u_{\{\LL\}} (S)$,  with $f\sim^u_{\{\LL\}}\sum_{j=0}^\infty c_j z^j$. Then, $\mathcal{T}_{\M} (f)\in \widetilde{\mathcal{A}}^u_{\{\LL\M\}} (S)$  with
%        \begin{equation}\label{eq.CoeffT_Mf}
%        	\mathcal{T}_{\M} (f)\sim^u_{\{\LL\M\}} \sum_{j=0}^{\infty}R_jc_j z^j,\qquad z\in S.
%        \end{equation}
%        Moreover, for any $A>0$, $\mathcal{T}_{\M} : \widetilde{\mathcal{A}}^u_{\LL, A} (S)\rightarrow\widetilde{\mathcal{A}}^u_{\LL\M, A } (S)$ is a continuous linear operator.
%\end{theorem}
%%
%\demo{Proof}
%By definition of $\widetilde{\mathcal{A}}^u_{\{\LL\}} (S)$ we have that $f$ is bounded in $S$ by some constant $C>0$. Since $\M$ is log-convex, we have that $M_j\le m_j^j$ for all $j\in\NN_0$ and then
%$$\sum_{j=0}^{\infty}\frac{1}{2^j} \frac{M_j}{m_j^j} \left|f(m_j z) \right| \leq  C  \sum_{j=0}^{\infty}\frac{1}{2^j}=2C, \qquad z\in S.$$
%Consequently, the series defining $\mathcal{T}_{\M} (f)$ normally converges in the whole of $S$, it provides a function being holomorphic in $S$.
%
%For each $z\in S$ and every $n\in\NN_0$, we observe then that
%\begin{align*}
%	&\left|\mathcal{T}_{\M} (f) (z) -  \sum_{k=0}^{n-1} R_kc_k z^k\right|= \left|\sum_{j=0}^{\infty}\frac{1}{2^j} \frac{M_j}{m_j^j} f(m_j z) -  \sum_{k=0}^{n-1} c_k  (\sum^{\infty}_{j=0} \frac{1}{2^j} \frac{M_j}{m_j^j}m_j^k)  z^k\right|
%	\\&
%	=\left|\sum_{j=0}^{\infty}\frac{1}{2^j} \frac{M_j}{m_j^j}  \left(f(m_j z) -  \sum_{k=0}^{n-1} c_k   (m_jz)^k\right)\right|
%	\\&
%	\le\sum_{j=0}^{\infty}\frac{1}{2^j}\frac{M_j}{m_j^j}C A^n L_n m_j^n |z|^n= C A^n L_n R_n |z|^n
%	\\&
%	\le 2 C A^n L_n M_n |z|^n,
%\end{align*}
%where the last estimate holds by Lemma \ref{Scomplemma}.
%Finally, suppose $f\in \widetilde{\mathcal{A}}^u_{\LL,A} (S)$ for some $A>0$, then for all $j\in\NN_0$ we can estimate
%\begin{align*}
%	&\left|\mathcal{T}_{\M} (f) (z) -  \sum_{k=0}^{n-1} R_kc_k z^k\right|\leq \sum_{j=0}^{\infty}\frac{1}{2^j} \frac{M_j}{m_j^j}  \left|f(m_j z) -  \sum_{k=0}^{n-1} c_k   (m_jz)^k\right|
%	\\&
%	\le \|f\|_{\M,A,\overset{\sim}{u}} A^n L_n|z|^n \sum_{j=0}^{\infty}\frac{1}{2^j}M_jm_j^{n-j}=\|f\|_{\M,A,\overset{\sim}{u}} A^n L_n R_n|z|^n.
%\end{align*}
%By Lemma~\ref{Scomplemma} we know that $R_n\le 2M_n$, so $\mathcal{T}_{\M} (f)\in\widetilde{\mathcal{A}}^u_{\LL\M, A} (S)$, and moreover
%$$\|\mathcal{T}_{\M} (f)\|_{\LL\M, A,\overset{\sim}{u}}=\sup_{z\in S,n\in\NN_{0}}\frac{|\mathcal{T}_{\M} (f) (z) -  \sum_{k=0}^{n-1} R_kc_k z^k|}{A^{n}L_{n}M_{n}|z|^n}\le 2\|f\|_{\M,A,\overset{\sim}{u}}.$$
%It follows that $\mathcal{T}_{\M} : \widetilde{\mathcal{A}}^u_{\LL, A} (S)\rightarrow\widetilde{\mathcal{A}}^u_{\LL\M, A } (S) $ is a well-defined continuous linear operator for any $A>0$.
%\qed\enddemo
%
%\begin{theorem}\label{CharacteristicTransformUniformAsymptotic1}
%	%Let $\M$ be a (lc) sequence, $\LL\in\RR_{>0}^{\NN_0}$ and for each $\beta>0$ take
% Let $\M$ be a (lc) sequence, $\LL\in\RR_{>0}^{\NN_0}$ and for a given sector $S$ take $f\in \widetilde{\mathcal{A}}^u_{\{\LL\}} (S)$  with $f\sim^u_{\{\LL\}}\sum_{j=0}^{\infty}c_j z^j$. If $(|c_j|)_{j\in\NN_0}$ is equivalent to $\LL$, then $(|c_j R_j|)_{j\in\NN_0}$ is
% equivalent to $\LL\M$. Consequently, $\mathcal{T}_{\M}(f)$ is characteristic in the class $\widetilde{\mathcal{A}}^u_{\{\LL\M\}} (S)$.
%\end{theorem}
%
%\demo{Proof}
%The first assertion is clear from Lemma~\ref{Scomplemma} and~\eqref{eq.CoeffT_Mf}. The second one stems from Theorems~\ref{CharacteristicTransformUniformAsymptotic} and~\ref{optthm1}.
%\qed\enddemo
%
%\subsection{Construction of characteristic functions}\label{ssect.ConstrCharactFunct}
%Given a sequence $\M\in\RR_{>0}^{\NN_0}$ and $\alpha>0$ we construct now, under suitable assumptions, characteristic functions in $\widetilde{\mathcal{A}}^u_{\{\M\}} (S_{\alpha})$.
%% and in $\widetilde{\mathcal{A}}^u_{\{\M\}} (S_{\alpha})$.
%For this we are using the basic functions from Subsection \ref{ssect.BasicFct} and the characteristic transform from Subsection \ref{ssect.CharactTransf}.\vspace{6pt}
%
%
%\begin{theorem}\label{mainTthm1}
%	Let $\M\in\RR_{>0}^{\NN_0}$ and $\alpha>0$.
%	%
%	\begin{enumerate}
%		\item If $\alpha< 2$, we assume that  $\overline{\G}^{2-\alpha}\M:=(j^{(2-\alpha)j}M_j)_{j\in\NN_0}$
%		is equivalent to an (lc) sequence $\LL$. Then, $\mathcal{T}_{\LL} (\widetilde{E}_{\alpha}) $ is  characteristic in the class $\widetilde{\mathcal{A}}^u_{\{\M\}} (S_{\alpha})$.
%		\item If $\alpha= 2$, we assume that  $\M$
%		is equivalent to an (lc) sequence $\LL$. Then, $\mathcal{T}_{\LL} (K_{RS}) $ is characteristic in the class $\widetilde{\mathcal{A}}^u_{\{\M\}} (S_{2})$.
%		\item If $\alpha>2$, we assume that there exists $\a'>\a$ such that $\overline{\G}^{2-\alpha'}\M:=(j^{(2-\alpha')j}M_j)_{j\in\NN_0}$
%		is equivalent to an (lc) sequence $\LL$. Then, $\mathcal{T}_{\LL} (f_{\alpha,\alpha'}) $ is characteristic in the class $\widetilde{\mathcal{A}}^u_{\{\M\}} (S_{\alpha})$.
%	\end{enumerate}
%\end{theorem}
%%
%%
%%
%\demo{Proof}
%This follows by Theorems~\ref{Ealphathm}, \ref{fgalphaalphaprimethm},  \ref{CharacteristicTransformUniformAsymptotic} and \ref{CharacteristicTransformUniformAsymptotic1}, and  from the fact that $\overline{\G}^{\alpha-2}\LL$ in case 1, resp. $\LL$ in case 2, resp. $\overline{\G}^{\alpha'-2}\LL$ in case 3 is equivalent to $\M$.
%\qed\enddemo
%
%\begin{remark}%\label{rem.gammaindexcharactfunct}
%	In order to guarantee that the hypotheses in the previous theorem are satisfied, one can compute the index $\gamma(\M)$ and check whether it is greater than $\a-2$. If this is the case, the very definition of this index implies that for {any $\beta$ such that $\gamma(\M)>\beta>\a-2$ the property $\left(P_{\beta}\right)$} (see Subsection~\ref{ssect.IndexgammaM}) is satisfied, and so there exists a suitable  (lc) sequence $\LL$ in the desired conditions.
%\end{remark}
%
%%%%%%%%%%%%%%%%%%%%%%%%%%%%%%%%%%%%%%%%%%%%%%%%%%%%%%%%%%%%%%%%%%%%%%%%%%%
%%%%%%%%%%%%%%%%%%%%%%%%%%%%%%%%%%%%%%%%%%%%%%%%%%%%%%%%%%%%%%%%%%%%%%%%%%%
\section{Characteristic functions for classes with Roumieu uniform asymptotics}\label{sect.CharactFunctions}


We introduce the concept of characteristic functions for the classes under consideration.


\begin{definition}\label{chardef}
	Let $\LL\in\RR_{>0}^{\NN_0}$ and $S$ be a given sector.
A function $f\in \widetilde{\mathcal{A}}^u_{\{\LL\}} (S)$ is said to be \textit{characteristic} in the class $\widetilde{\mathcal{A}}^u_{\{\LL\}} (S)$ if, whenever $f\in \widetilde{\mathcal{A}}^u_{\{\M\}} (S)\subseteq \widetilde{\mathcal{A}}^u_{\{\LL\}} (S)$ for some $\M\in\RR_{>0}^{\NN_0}$, we have that $\widetilde{\mathcal{A}}^u_{\{\M\}} (S)= \widetilde{\mathcal{A}}^u_{\{\LL\}} (S)$.
\end{definition}

% Let us fix some notation. If $f\in \widetilde{\mathcal{A}}^u_{\{\LL\}} (S)$  with $f\sim^u_{\{\LL\}} \sum_{j=0}^\infty c_j z^j$, then put

% %$$r_{f}(z,n):= \frac{f(z)-\sum_{j=0}^{n-1} c_j z^j}{z^n},\;\;\;z\in S,\;n\in\NN_0.$$
% $$r_{f}(z,n):= f(z)-\sum_{j=0}^{n-1} c_j z^j,\;\;\;z\in S,\;n\in\NN_0.$$

% We observe that $z r_f(z,n+1)/z^{n+1}=r_f(z,n)/z^{n}-c_n$. If $f\in\widetilde{\mathcal{A}}^u_{\{\LL\}} (S)$, then $z\mapsto z^{-n}r_f(z,n)$ is bounded in $S$ for each $n\in\NN_0$ fixed and so
% $$\forall\;n\in\NN_0:\;\;\;\lim_{z\in S,z\rightarrow 0} z^{-n}r_f(z,n)=c_n.$$
Let $f\in \widetilde{\mathcal{A}}^u_{\{\LL\}} (S)$. We consider the sequence defined by
\begin{equation*}%\label{tildesequenceC}
	\widetilde{C}_n(f):=\sup_{z\in S}|z^{-n}r_f(z,n)|, \qquad n\in\NN_0.
\end{equation*}

We deduce that a function is characteristic of the class if the sequence of moduli of the coefficients of the associated asymptotic expansion (resp. the sequence $(\widetilde{C}_j(f))_{j\in\NN_0}$) is equivalent to the sequence defining the class. More precisely,

\begin{theorem}\label{optthm1}
	Let $\LL\in\RR_{>0}^{\NN_0}$, $S$ be a given sector and $f\in \widetilde{\mathcal{A}}^u_{\{\LL\}} (S)$ with $f\sim^u_{\{\LL\}} \sum_{j=0}^\infty c_j z^j$. Then, each of the following  conditions implies the next one:
	\begin{enumerate}
		\item[$(1)$] The sequence $(|c_j|)_{j\in\NN_0}$ is equivalent to $\LL$.
		\item[$(2)$] The sequence $(\widetilde{C}_j(f))_{j\in\NN_0}$ is equivalent to $\LL$.
		\item[$(3)$] $f$ is characteristic in the class $\widetilde{\mathcal{A}}^u_{\{\LL\}} (S)$.
	\end{enumerate}
\end{theorem}

\demo{Proof}
(1) $\Rightarrow$ (2) As $f\in \widetilde{\mathcal{A}}^u_{\{\LL\}} (S)$, there exist $A,B>0$ such that $\widetilde{C}_n(f)\le AB^nL_n$ for every $n\in\NN_0$. On the other hand, it is clear that $\widetilde{C}_n(f)\geq |c_n|$, and the hypothesis allows us to conclude the other estimate.
%The estimate from above for $\widetilde{C}_n(f)$ follows by definition and the estimate from below holds since $\widetilde{C}_n(f)\geq |c_n|$.

(2)  $\Rightarrow$ (3) By assumption, there exist $A,B>0$ such that $L_n\le A{B}^n\widetilde{C}_n(f)$ for every $n\in\NN_0$. If for some $\M=(M_n)_{n\in\NN_0}\in\RR_{>0}^{\NN_0}$ we have $f\in \widetilde{\mathcal{A}}^u_{\{\M\}} (S)\subseteq \widetilde{\mathcal{A}}^u_{\{\LL\}} (S)$, there exist $C,D>0$ such that $\widetilde{C}_n(f)\le C {D}^nM_n$ for every $n\in\NN_0$. The two deduced inequalities show that $L_n\le{ AC(BD)^n}M_n$ for every $n\in\NN_0$, what easily implies that $\widetilde{\mathcal{A}}^u_{\{\LL\}} (S)\subseteq \widetilde{\mathcal{A}}^u_{\{\M\}} (S)$, and we are done.
\qed\enddemo

%For $f\in \mathcal{A}_{\{\LL\}}(S)$ we consider the sequence defined by
%\begin{equation*}%\label{sequenceC}
%	C_n(f):=\sup_{z\in S} |f^{(n)}(z)|, \qquad n\in\NN_0.
%\end{equation*}
%The next statement provides conditions on $f$ which imply it is characteristic.
%%Similarly as before we get:
%
%\begin{theorem}\label{optthm2}
%	Let $\LL\in\RR_{>0}^{\NN_0}$, $S$ be a given sector and $f\in \mathcal{A}_{\{\LL\}} (S)$. Then, each of the following  conditions implies the next one:
%	\begin{enumerate}
%		\item[$(1)$] The sequence $(|f^{(j)}(0)|)_{j\in\NN_0}$ is equivalent to $\LL$.
%		\item[$(2)$] The sequence $(C_j(f))_{j\in\NN_0}$ is equivalent to $\LL$.
%		\item[$(3)$] $f$ is characteristic in the class $\mathcal{A}_{\{\LL\}} (S)$.
%	\end{enumerate}
%\end{theorem}
%
%\demo{Proof}
%(1) $\Rightarrow$ (2) As $f\in \mathcal{A}_{\{\LL\}} (S)$, there exist $A,B>0$ such that $C_n(f)\le AB^nL_n$ for every $n\in\NN_0$. On the other hand, it is clear that $C_n(f)\geq |f^{(n)}(0)|$, and the hypothesis allows us to conclude the other estimate.
%
%(2)  $\Rightarrow$ (3) By assumption, there exist $A,B>0$ such that $L_n\le A{B}^nC_n(f)$ for every $n\in\NN_0$. If for some $\M=(M_n)_{n\in\NN_0}\in\RR_{>0}^{\NN_0}$ we have $f\in \mathcal{A}_{\{\M\}} (S)\subseteq \mathcal{A}_{\{\LL\}} (S)$, there exist $C,D>0$ such that $C_n(f)\le C {D}^nM_n$ for every $n\in\NN_0$. The two deduced inequalities show that $L_n\le{ AC(BD)^n}M_n$ for every $n\in\NN_0$, what easily implies that $\mathcal{A}_{\{\LL\}} (S)\subseteq \mathcal{A}_{\{\M\}} (S)$, and we are done.
%\qed\enddemo

\subsection{Basic functions}\label{ssect.BasicFct}

Recall the notations {$\G^s:=(j!^s)_{j\in\NN_0}$ and }$\overline{\G}^s:=(j^{js})_{j\in\NN_0}$, $s\in\RR$, and that $\overline{\G}^s\hyperlink{approx}{\approx}\G^s$, see \eqref{Stirling}. We are going to construct characteristic functions in the uniform asymptotic classes associated with $\overline{\G}^s$, where $s$ depends on the opening of the sector.%\vspace{6pt}

%In fact, we need in the following $s\in\RR$ (so also negative $s$ are allowed) but

The \textit{{two-parametric} Mittag-Leffler function} is defined for all {complex parameters $A, B$ with $\Re(A)>0$} by
$$E_{A,B}(z):=\sum_{j=0}^\infty \frac{z^j}{\Gamma (Aj+B)},\quad z\in\CC,$$
where $\Gamma$ denotes the Gamma function.
For the construction of characteristic functions in sectors $S_\alpha$ for $\alpha\in(0,1]$ we will take $A=2-\alpha$ and $B=4-\alpha$ and we set
\begin{equation*}%\label{goodEalpha}
	\widetilde{E}_{\alpha}(z):=E_{2-\alpha, 4-\alpha}(-z){=\sum_{j=0}^\infty \frac{(-1)^jz^j}{\Gamma ((2-\alpha)(j+1)+2)},\quad z\in\CC}.
\end{equation*}
%For the construction of characteristic functions in sectors $S_\alpha$ for $\alpha\in[0,2)$ we will take $A=2-\alpha$ and $B=4-\alpha$ and so we set
%\begin{equation}\label{goodEalpha}
%	\widetilde{E}_{\alpha}(z):=E_{2-\alpha, 4-\alpha}(-z).
%\end{equation}

We recall the following statement.

\begin{theorem}\label{Ealphathm} (\cite[Thm. 5, Thm. 20]{salinas62})
	Let $\alpha\in(0,2)$, then
	\begin{equation*}%\label{Ealphagoodestimate}
		\exists\;C\ge 1\;\forall\;z\in S_{\alpha}\;\forall\;n\in\NN_0:\;\;\;\left| \widetilde{E}_{\alpha}(z)-\sum_{j=0}^{n-1} \frac{(-z)^j}{\Gamma ((j+1) (2-\alpha)+2)}\right|\leq C \frac{e^n}{n^{(2-\alpha)n}}|z|^n.
	\end{equation*}
	Consequently, $\widetilde{E}_{\alpha}\in\widetilde{\mathcal{A}}^u_{\{\overline{\G}^{\alpha-2}\}}(S_\alpha)$ and,  
    %$\widetilde{E}_{\alpha}\in \mathcal{A}_{\{\overline{\G}^{\alpha-1}\}}(S_\alpha)$.	
    moreover, 
    $\widetilde{E}_{\alpha}$ is a characteristic function in the class $\widetilde{\mathcal{A}}^u_{\{\overline{\G}^{\alpha-2}\}}(S_\alpha)$.
%	Let $\widetilde{E}_{\alpha}(z)$ be the function from \eqref{goodEalpha}.
%\red{CASE $\alpha=0$.}
%	\begin{itemize}
%		\item[$(i)$] By \cite[Thm. 5]{salinas62} we get: Let $\alpha\in[0,2)$, then
%		$$\exists\;C\ge 1\;\forall\;z\in S_{\alpha}\;\forall\;n\in\NN_0:\;\;\;\left| \widetilde{E}_{\alpha}(z)-\sum_{j=0}^{n-1} \frac{(-z)^j}{\Gamma ((j+1) (2-\alpha)+2)}\right|\leq C \frac{e^n}{n^{(2-\alpha)n}}|z|^n.$$
%		Consequently, $\widetilde{E}_{\alpha}(z)\in\widetilde{\mathcal{A}}^u_{\{\overline{\G}^{\alpha-2}\}}(S_\alpha)$ and, moreover, $\widetilde{E}_{\alpha}(z)$ is a characteristic function in the class $\widetilde{\mathcal{A}}^u_{\{\overline{\G}^{\alpha-2}\}}(S_\alpha)$.
%	\end{itemize}
%		\item[$(ii)$]
%Let $\alpha\in(0,1]$, then
%		\begin{equation}\label{Ealphagoodestimate}
%			\forall\;z\in S_{\alpha}\;\forall\;n\in\NN_0:\;\;\;\left|\widetilde{E}_{\alpha}^{(n)}(z)\right|\leq 2 \frac{n! e^n}{n^{(2-\alpha)n}}.
%		\end{equation}
%		Consequently, $\widetilde{E}_{\alpha}\in \mathcal{A}_{\{\overline{\G}^{\alpha-1}\}}(S_\alpha)$.
%		Moreover,
%{$$
%\widetilde{E}_{\alpha}^{(n)}(0)=\frac{(-1)^n n!}{\Gamma ((2-\alpha)(n+1)+2)},\quad n\in\NN_0,
%$$
%and so} $\widetilde{E}_{\alpha}$ is a characteristic function in the class $\mathcal{A}_{\{\overline{\G}^{\alpha-1}\}}(S_\alpha)$.
%	\end{itemize}
\end{theorem}

Next we consider the Rodr\'{i}guez-Salinas function (see \cite{salinas62})
\begin{equation*}%\label{RSfunction}
	K_{RS}(z):=\frac{(1+z)\log(1+z)-z}{z^2},
\end{equation*}
which is clearly holomorphic in $S_2$.

This function satisfies the following theorem %Now we recall \cite[Thm. 3]{salinas62}

\begin{theorem}\label{RStheorem} (\cite[Thm. 3]{salinas62})
	There exists a constant $C$ (we can take $C=4$) such that
	 $$\forall\;z\in S_2%\CC\backslash(-\infty,-1]
	 \;\forall\;n\in\NN_0:\;\;\;\left| K_{RS}(z)-\sum_{j=0}^{n-1} \frac{(-z)^j}{(j+1)(j+2)}\right|\leq C|z|^n.$$
	Consequently, for the constant sequence $\mathbf{1}:=(1)_{j\in\NN_0}$ we have that  $K_{RS}\in\widetilde{\mathcal{A}}^u_{\{\mathbf{1}\}}(S_2)$ and, moreover, $K_{RS}$ is a characteristic function in the class $\widetilde{\mathcal{A}}^u_{\{ \mathbf{1}\}}(S_2)$.
\end{theorem}
%
%
Finally, let $\alpha>2$ and take $\alpha'>\alpha$. For all $z\in S_\alpha$ we define
\begin{equation}\label{fctfalphaalphaprime}
	f_{\alpha,\alpha'}(z):= \int_{0}^{\infty(-\phi)} K_{RS}(zv^{\alpha'-2}) e^{-v} dv,
\end{equation}
where we choose $\phi  \in (-\frac{(\alpha-2)}{(\alpha'-2)}\frac{\pi}{2}, \frac{(\alpha-2)}{(\alpha'-2)}\frac{\pi}{2})$ with $|\arg(z)- (\alpha'-2)\phi|<\pi$.\\
\begin{theorem}\label{fgalphaalphaprimethm} (\cite[Thm. 4, Thm. 8]{salinas62})
	Let $\alpha>2$, $\alpha'>\alpha$ and
	$f_{\alpha,\alpha'}$ be the function from
	\eqref{fctfalphaalphaprime}. Then
				\begin{align*}%\label{falphagoodestimate}
						\exists\;C,A\ge 1\;\forall\;z\in S_{\alpha}\;\forall\;n\in\NN_0:\;\;\;&\left|f_{\alpha,\alpha'} (z) -\sum_{j=0}^{n-1} \frac{(-z)^j}{(j+1)(j+2)} \Gamma (  (\alpha'-2)j+1)\right|\nonumber
						\\&
						\le CA^n \Gamma ( (\alpha'-2)n+1)|z|^n.
					\end{align*}
				Consequently, $f_{\alpha,\alpha'}\in\widetilde{\mathcal{A}}^u_{\{\overline{\G}^{\alpha'-2}\}}(S_\alpha)$ and,
				moreover, $f_{\alpha,\alpha'}$ is a characteristic function in the class $\widetilde{\mathcal{A}}^u_{\{ \overline{\G}^{\alpha'-2} \}}(S_\alpha)$.
	\end{theorem}
%Let $\alpha>1$ and take $\alpha'>\alpha$. For all $z\in S_\alpha$ we define
%\begin{equation}\label{fctgalphaalphaprime}
%	g_{\alpha,\alpha'} (z):= \int_{0}^{\infty(-\phi)} e^{-zv^{\alpha'-1}} e^{-v} dv,
%\end{equation}
%where we choose $\phi  \in (-\frac{(\alpha-1)}{(\alpha'-1)}\frac{\pi}{2}, \frac{(\alpha-1)}{(\alpha'-1)}\frac{\pi}{2})$ with $|\arg(z)- (\alpha'-1)\phi|<\pi/2$.
%
%\begin{theorem}\label{fgalphaalphaprimethm} (\cite[Thm. 28]{salinas62})
%	Let $\alpha>1$, $\alpha'>\alpha$ and
%%$f_{\alpha,\alpha'}$ resp.
%$g_{\alpha,\alpha'}$ be the function from
%%\eqref{fctfalphaalphaprime} resp.
%\eqref{fctgalphaalphaprime}.
%%	\begin{itemize}
%%		\item[$(i)$] By \cite[Thm. 4, Thm. 8]{salinas62} we get: Let $\alpha>2$ and take $\alpha'>\alpha$. Then
%%		\begin{align*}
%%			\exists\;C,A\ge 1\;\forall\;z\in S_{\alpha}\;\forall\;n\in\NN_0:\;\;\;&\left|f_{\alpha,\alpha'} (z) -\sum_{j=0}^{n-1} \frac{(-z)^j}{(j+1)(j+2)} \Gamma (  (\alpha'-2)j+1)\right|
%%			\\&
%%			\le CA^n \Gamma ( (\alpha'-2)n+1)|z|^n.
%%		\end{align*}
%%		Consequently, we have that $f_{\alpha,\alpha'}\in\widetilde{\mathcal{A}}^u_{\{\overline{\G}^{\alpha'-2}\}}(S_\alpha)$ and,
%%		moreover, $f_{\alpha,\alpha'}$ is a characteristic function in the class $\widetilde{\mathcal{A}}^u_{\{ \overline{\G}^{\alpha'-2} \}}(S_\alpha)$.
%		
%%		\item[$(ii)$]
%Then,
%		\begin{equation}\label{galphagoodestimate}
%			\exists\;C,A\ge 1\;\forall\;z\in S_{\alpha}\;\forall\;n\in\NN_0:\;\;\;\left| g_{\alpha,\alpha'}^{(n)}(z)\right|\leq C A^n \Gamma ((\alpha'-1)n+1).
%		\end{equation}
%		Consequently, $g_{\alpha,\alpha'}\in \mathcal{A}_{\{\overline{\G}^{\alpha'-1}\}}(S_\alpha)$. Moreover,
%{$$
%g_{\alpha,\alpha'}^{(n)}(0)=(-1)^n  \Gamma((\alpha'-1)n+1),\qquad n\in\NN_0,
%$$
%and so} $g_{\alpha,\alpha'}$ is a characteristic function of the class $\mathcal{A}_{\{\overline{\G}^{\alpha'-1}\}}(S_\alpha)$.
%%\end{itemize}
%\end{theorem}

\subsection{Characteristic transform}\label{ssect.CharactTransf}
% The $\mathcal{T}_{\M}$ transform originally appeared in the work of Rodríguez Salinas~\cite{salinas62}, allowed us in a previous work by the authors \cite{ultraholomstab}, the construction of characterisitc functions with a precise control of the derivatives at 0 in ultraholomorphic classes defined by uniform bounds on the derivatives. Now, the same functional transform $\mathcal{T}_{\M}$ modifies the coefficients of the asymptotic expansion of a function $f\in\widetilde{\mathcal{A}}^u_{\{\LL\}} (S)$ in order to obtain a new function in the class $\widetilde{\mathcal{A}}^u_{\{\LL\M\}} (S)$, under suitable hypotheses. This transformation allows us to construct characteristic functions in more general classes than the Gevrey ones considered previously.
% Following again the work of Rodríguez Salinas~\cite{salinas62}, we present a functional transform $\mathcal{T}_{\M}$ that modifies the coefficients of the asymptotic expansion of a function $f\in\widetilde{\mathcal{A}}^u_{\{\LL\}} (S)$ in order to obtain a new function in the class $\widetilde{\mathcal{A}}^u_{\{\LL\M\}} (S)$, under suitable hypotheses. This transformation allows us to construct characteristic functions in more general classes than the Gevrey ones considered previously.
% 
% 
The $\mathcal{T}_{\M}$ transform originally appeared in the work of Rodríguez Salinas~\cite{salinas62} and was used by the authors in the previous work~\cite{ultraholomstab} to construct characteristic functions with precise control of the derivatives at $0$ in ultraholomorphic classes defined by uniform bounds on the derivatives.

In this context, the same functional transform $\mathcal{T}_{\M}$ modifies the coefficients of the asymptotic expansion of a function $f\in\widetilde{\mathcal{A}}^u_{\{\LL\}} (S)$ to obtain a new function in the class $\widetilde{\mathcal{A}}^u_{\{\LL\M\}} (S)$, under suitable hypotheses. This transformation allows us to construct characteristic functions in classes more general than the Gevrey classes considered previously.
\begin{definition}
	Let $\M$ be an \hyperlink{lc}{$(\text{lc})$} sequence, $\LL\in\RR_{>0}^{\NN_0}$, $S$ a sector and $f\in\widetilde{\mathcal{A}}^u_{\{\LL\}} (S)$. 
	%resp. $f\in \mathcal{A}_{\{\LL\}} (S)$
Then we define the $\mathcal{T}_{\M}-$transform of $f$ by
	\begin{equation*}%\label{Ttrasform}
		\mathcal{T}_{\M}(f)(z):= \sum_{j=0}^{\infty}\frac{1}{2^j} \frac{M_j}{m_j^j} f(m_j z), \qquad z\in S.
	\end{equation*}
\end{definition}
This expression should be compared with the characteristic functions obtained in the ultradifferentiable and ultraholomorphic setting in \cite[Thm. 1]{thilliez}, \cite[Lemma 2.9]{compositionpaper}, \cite[Lemma 1]{Poschel} and \cite[Definition 4.5.]{ultraholomstab}. 
For every $j\in\NN_0$ let us set
\begin{equation*}%\label{Sexpression}
	R_j:=\sum^{\infty}_{n=0} \frac{1}{2^n} \frac{M_n}{m_n^n}m_n^j.
\end{equation*}
The following result provides estimates for this sequence in terms of the general sequence $\M$ we depart from.

\begin{lemma}\label{Scomplemma}(\cite[Lemma 4.6]{ultraholomstab})
	Let $\M\in\RR_{>0}^{\NN_0}$, then
	$$\forall\;j\in\NN_0:\;\;\;R_j\ge\frac{1}{2^j} M_j.$$
	If $\M$ is \hyperlink{lc}{$(\text{lc})$}, then also
	$$\forall\;j\in\NN_0:\;\;\;R_j\le 2 M_j,$$
	and so $(R_j)_{j\in\NN_0}$ is equivalent to $\M$.
\end{lemma}

% \demo{Proof}
% For any $j\in\NN_0$ we choose $n=j$ in the sum and get $R_j\ge\frac{1}{2^j} \frac{M_j}{m_j^j}m_j^j=\frac{1}{2^j} M_j$.

% For the converse we recall that since $\M$ is (lc) we have $m_0\le m_1\le\dots$ and so
% $$\forall\;j,n\in\NN_0:\;\;\;(m_n)^{j-n}\le\frac{M_j}{M_n},$$
% see \cite[Thm. 1]{thilliez} and the detailed proof in \cite[(3.1.2)]{diploma}. Thus
% $$R_j=\sum^{\infty}_{n=0}\frac{1}{2^n}M_n m_n^{j-n}\le\sum^{\infty}_{n=0}\frac{1}{2^n}M_n\frac{M_j}{M_n}=2M_j$$
% for all $j\in\NN_0$.
% \qed\enddemo

\begin{theorem}\label{CharacteristicTransformUniformAsymptotic}
	%Let $\M$ be a (lc) sequence, $\LL\in\RR_{>0}^{\NN_0}$ and for each $\beta>0$ take
        Let $\M$ be a \hyperlink{lc}{$(\text{lc})$} sequence, $\LL\in\RR_{>0}^{\NN_0}$ and for a given sector $S$ take $f\in\widetilde{\mathcal{A}}^u_{\{\LL\}} (S)$,  with $f\sim^u_{\{\LL\}}\sum_{j=0}^\infty c_j z^j$. Then, $\mathcal{T}_{\M} (f)\in \widetilde{\mathcal{A}}^u_{\{\LL\M\}} (S)$  with 
        \begin{equation}\label{eq.CoeffT_Mf}
        	\mathcal{T}_{\M} (f)\sim^u_{\{\LL\M\}} \sum_{j=0}^{\infty}R_jc_j z^j,\qquad z\in S.
        \end{equation}
        Moreover, for any $A>0$, $\mathcal{T}_{\M} : \widetilde{\mathcal{A}}^u_{\LL, A} (S)\rightarrow\widetilde{\mathcal{A}}^u_{\LL\M, A } (S)$ is a continuous linear operator.
\end{theorem}
%
\demo{Proof}
By definition of $\widetilde{\mathcal{A}}^u_{\{\LL\}} (S)$ we have that $f$ is bounded in $S$ by some constant $C>0$. Since $\M$ is log-convex, we have that $M_j\le m_j^j$ for all $j\in\NN_0$ and then
$$\sum_{j=0}^{\infty}\frac{1}{2^j} \frac{M_j}{m_j^j} \left|f(m_j z) \right| \leq  C  \sum_{j=0}^{\infty}\frac{1}{2^j}=2C, \qquad z\in S.$$
Consequently, the series defining $\mathcal{T}_{\M} (f)$ normally converges in the whole of $S$, it provides a function holomorphic in $S$.
%and so the proof follows as in Theorem \ref{CharacteristicTransformUniformAsymptotic}.
%	\item
For each $z\in S$ and every $n\in\NN_0$, we observe then that
\begin{align*}
	\left|\mathcal{T}_{\M} (f) (z) -  \sum_{k=0}^{n-1} R_kc_k z^k\right|&= \left|\sum_{j=0}^{\infty}\frac{1}{2^j} \frac{M_j}{m_j^j} f(m_j z) -  \sum_{k=0}^{n-1} c_k  (\sum^{\infty}_{j=0} \frac{1}{2^j} \frac{M_j}{m_j^j}m_j^k)  z^k\right|
	\\&
	=\left|\sum_{j=0}^{\infty}\frac{1}{2^j} \frac{M_j}{m_j^j}  \left(f(m_j z) -  \sum_{k=0}^{n-1} c_k   (m_jz)^k\right)\right|
	\\&
	\le\sum_{j=0}^{\infty}\frac{1}{2^j}\frac{M_j}{m_j^j}C A^n L_n m_j^n |z|^n= C A^n L_n R_n |z|^n
	\\&
	\le 2 C A^n L_n M_n |z|^n,
\end{align*}
where the last estimate holds by Lemma \ref{Scomplemma}.
Finally, suppose $f\in \widetilde{\mathcal{A}}^u_{\LL,A} (S)$ for some $A>0$, then for all $n\in\NN_0$ we can estimate
\begin{align*}
	&\left|\mathcal{T}_{\M} (f) (z) -  \sum_{k=0}^{n-1} R_kc_k z^k\right|\leq \sum_{j=0}^{\infty}\frac{1}{2^j} \frac{M_j}{m_j^j}  \left|f(m_j z) -  \sum_{k=0}^{n-1} c_k   (m_jz)^k\right|
	\\&
	\le \|f\|_{\M,A,\overset{\sim}{u}} A^n L_n|z|^n \sum_{j=0}^{\infty}\frac{1}{2^j}M_jm_j^{n-j}=\|f\|_{\M,A,\overset{\sim}{u}} A^n L_n R_n|z|^n.
\end{align*}
By Lemma~\ref{Scomplemma} we know that $R_n\le 2M_n$, so $\mathcal{T}_{\M} (f)\in \mathcal{A}_{\LL\M, A} (S)$, and moreover
$$\|\mathcal{T}_{\M} (f)\|_{\LL\M, A,\overset{\sim}{u}}=\sup_{z\in S,n\in\NN_{0}}\frac{|\mathcal{T}_{\M} (f) (z) -  \sum_{k=0}^{n-1} R_kc_k z^k|}{A^{n}L_{n}M_{n}|z|^n}\le 2\|f\|_{\M,A,\overset{\sim}{u}}.$$
It follows that $\mathcal{T}_{\M} : \widetilde{\mathcal{A}}^u_{\LL, A} (S)\rightarrow\widetilde{\mathcal{A}}^u_{\LL\M, A } (S) $ is a well-defined continuous linear operator for any $A>0$.
\qed\enddemo

\begin{theorem}\label{CharacteristicTransformUniformAsymptotic1}
	%Let $\M$ be a (lc) sequence, $\LL\in\RR_{>0}^{\NN_0}$ and for each $\beta>0$ take
 Let $\M$ be a \hyperlink{lc}{$(\text{lc})$} sequence, $\LL\in\RR_{>0}^{\NN_0}$ and for a given sector $S$ take $f\in \widetilde{\mathcal{A}}^u_{\{\LL\}} (S)$  with $f\sim^u_{\{\LL\}}\sum_{j=0}^{\infty}c_j z^j$. If $(|c_j|)_{j\in\NN_0}$ is equivalent to $\LL$, then $(|c_j R_j|)_{j\in\NN_0}$ is
 equivalent to $\LL\M$. Consequently, $\mathcal{T}_{\M}(f)$ is characteristic in the class $\widetilde{\mathcal{A}}^u_{\{\LL\M\}} (S)$.
\end{theorem}

\demo{Proof}
The first assertion is clear from Lemma~\ref{Scomplemma}. The second one stems from Theorems~\ref{CharacteristicTransformUniformAsymptotic} and~\ref{optthm1}.
%By Theorem \ref{CharacteristicTransformUniformAsymptotic} we obtain $\mathcal{T}_{\M} (f)\sim^u_{\{\LL\M\}} \sum_{j=0}^{\infty}c_j R_j z^j$ in $S$. Since $(|c_j|)_{j\in\NN_0}$ is equivalent to $\LL$ and by Lemma \ref{Scomplemma} we get that
%$$\exists\;C,A\le 1\;\forall\;n\in\NN_0:\;\;\;|c_n R_n |\ge C A^n L_n M_n,$$
%and the converse estimate follows analogously. Thus $(|c_j R_j|)_{j\in\NN_0}$ is equivalent to $\LL\M$.
\qed\enddemo

%Note that for any $f\in\widetilde{\mathcal{A}}^u_{\{\LL\}} (S)$ we have
%$$\exists\;C,A\ge 1\;\forall\;n\in\NN_0:\;\;\;\widetilde{C}_n(f)\le CA^n L_n,$$
%with $\widetilde{C}_n$ denoting the value defined in \eqref{tildesequenceC}. So, in view of Theorem \ref{optthm1}, in order to see that if $f$ is characteristic we only need to check the bounds from below. Thanks to the equivalence of $(R_j)_{j\in\NN_0}$ and $\M$ verified in Lemma \ref{Scomplemma} we see that if $f\in \widetilde{\mathcal{A}}^u_{\{\LL\}} (S)$ is characteristic then $\mathcal{T}_{\M} (f)$ is characteristic in $\widetilde{\mathcal{A}}^u_{\{\LL\M\}}(S)$, too.

%\begin{theorem}\label{CharacteristicTransformUniformAsymptotic1}
%	%Let $\M$ be a (lc) sequence, $\LL\in\RR_{>0}^{\NN_0}$ and for each $\beta>0$ take
% Let $\M$ be a (lc) sequence, $\LL\in\RR_{>0}^{\NN_0}$ and for a given sector $S$ take $f\in \widetilde{\mathcal{A}}^u_{\{\LL\}} (S)$  with $f\sim^u_{\{\LL\}}\sum_{j=0}^{\infty}c_j z^j$. Assume that $(|c_j|)_{j\in\NN_0}$ is equivalent to $\LL$, then $\mathcal{T}_{\M} (f)\sim^u_{\{\LL\M\}} \sum_{j=0}^{\infty}c_j R_j z^j$ and $(|c_j R_j|)_{j\in\NN_0}$ is
%	equivalent to $\LL\M$.
%	
%	Consequently, $\mathcal{T}_{\M}(f)$ is characteristic in the class $\widetilde{\mathcal{A}}^u_{\{\LL\M\}} (S)$.
%\end{theorem}

%\demo{Proof}
%By Theorem \ref{CharacteristicTransformUniformAsymptotic} we obtain $\mathcal{T}_{\M} (f)\sim^u_{\{\LL\M\}} \sum_{j=0}^{\infty}c_j R_j z^j$ in $S$. Since $(|c_j|)_{j\in\NN_0}$ is equivalent to $\LL$ and by Lemma \ref{Scomplemma} we get that
%$$\exists\;C,A\le 1\;\forall\;n\in\NN_0:\;\;\;|c_n R_n |\ge C A^n L_n M_n,$$
%and the converse estimate follows analogously. Thus $(|c_j R_j|)_{j\in\NN_0}$ is equivalent to $\LL\M$.
%\qed\enddemo
%
%
%    THE CHARACTERISTIC TRANSFORM FOR UNIFORM ASYMPTOTIC CLASS
%
%
%
%
%
%\begin{theorem}\label{CharacteristicTransform}
%	%Let $\M$ be a (lc) sequence, $\LL\in\RR_{>0}^{\NN_0}$ and for each $\beta>0$ take
% Let $\M$ be a (lc) sequence, $\LL\in\RR_{>0}^{\NN_0}$ and for a given sector $S$ take $f\in \mathcal{A}_{\{\LL\}} (S)$.%
% %For all $n\in\NN_0$ put  $$b_n:=\lim_{z\in S, z\rightarrow 0} f^{(n)}(z),$$ and
% Then,
%%  we have:
%%	\begin{enumerate}
%%		\item $\mathcal{T}_{\M} (f)$ is holomorphic in $S$.
%%		\item For every $j\in\NN_0$ we define
%%		$$a_j:=(f^{(j)}(0)) R_j,$$
%%		with $R_j$ being the numbers from \eqref{Sexpression}. Then
%$\mathcal{T}_{\M} (f)\in \mathcal{A}_{\{\LL\M\}} (S)$ with
%\begin{equation}\label{eq.DerivAt0T_Mf}
%\mathcal{T}_{\M} (f)^{(j)}(0)=R_jf^{(j)}(0),\qquad j\in\NN_0.
%\end{equation}
%Moreover, for any $A>0$, $\mathcal{T}_{\M} : \mathcal{A} _{\LL, A} (S)\rightarrow\mathcal{A}_{\LL\M, A } (S)$ is a continuous linear operator.
%%	\end{enumerate}
%\end{theorem}
%
%\demo{Proof}
%%\begin{enumerate}
%%	\item
%By definition of $\mathcal{A}_{\{\LL\}} (S)$ we have that $f$ is bounded in $S$ by some constant $C>0$. Since $\M$ is log-convex, we have that $M_j\le m_j^j$ for all $j\in\NN_0$ and then
%	$$\sum_{j=0}^{\infty}\frac{1}{2^j} \frac{M_j}{m_j^j} \left|f(m_j z) \right| \leq  C  \sum_{j=0}^{\infty}\frac{1}{2^j}=2C, \qquad z\in S.$$
%Consequently, the series defining $\mathcal{T}_{\M} (f)$ normally converges in the whole of $S$, it provides a function holomorphic in $S$, and differentiation and limits can be interchanged with summation.
% %and so the proof follows as in Theorem \ref{CharacteristicTransformUniformAsymptotic}.
%%	\item
%For each $z\in S$ and every $j\in\NN_0$ we observe then that
%	$$(\mathcal{T}_{\M} (f))^{(j)} (z)= \sum_{n=0}^{\infty}\frac{1}{2^n} \frac{M_n}{m_n^n} m_n^jf^{(j)}(m_n z),$$
%and so
%$$
%\mathcal{T}_{\M} (f)^{(j)}(0)=\sum_{n=0}^{\infty}\frac{1}{2^n} \frac{M_n}{m_n^n} m_n^jf^{(j)}(0)=R_jf^{(j)}(0),\qquad j\in\NN_0,$$
%as desired.
%
%Suppose $f\in \mathcal{A}_{\LL,A} (S)$ for some $A>0$, then for all $j\in\NN_0$ we can estimate
%\begin{align*}
%|(\mathcal{T}_{\M} (f))^{(j)} (z)|&\le \sum_{n=0}^{\infty}\frac{1}{2^n} \frac{M_n}{m_n^n}m_n^j |f^{(j)}(m_nz)|
% \\&
% \le \|f\|_{\M,A} A^j L_j \sum_{n=0}^{\infty}\frac{1}{2^n}M_nm_n^{j-n}=\|f\|_{\M,A} A^j L_j R_j.
%\end{align*}
%By Lemma~\ref{Scomplemma} we know that $R_j\le 2M_j$, so $\mathcal{T}_{\M} (f)\in \mathcal{A}_{\LL\M, A } (S)$, and moreover
%$$\|\mathcal{T}_{\M} (f)\|_{\LL\M, A}=\sup_{z\in S}\frac{|(\mathcal{T}_{\M} (f))^{(j)} (z)|}{A^j L_j M_j}\le 2\|f\|_{\M,A}.$$
%It follows that $\mathcal{T}_{\M} : \mathcal{A} _{\LL, A} (S)\rightarrow\mathcal{A}_{\LL\M, A } (S)$ is a well-defined continuous linear operator for any $A>0$.
%%\end{enumerate}
%\qed\enddemo
%
%    THE CHARACTERISTIC TRANSFORM FOR ULTRAHOLOMORPHIC CLASS IS SUPPRESSED
%
%
%
%
%
%
%\begin{theorem}\label{CharacteristicTransform1}
%	%Let $\M$ be a (lc) sequence, $\LL\in\RR_{>0}^{\NN_0}$ and for each $\beta>0$ take
% Let $\M$ be a (lc) sequence, $\LL\in\RR_{>0}^{\NN_0}$ and for a given sector $S$ take $f\in \mathcal{A}_{\{\LL\}} (S)$. %
% %For all $n\in\NN_0$ we denote  $$b_n:=\lim_{z\in S, z\rightarrow 0} f^{(n)}(z).$$
%	If $(|f^{(j)}(0)|)_{j\in\NN_0}$ is equivalent to $\LL$, then
%$(|\mathcal{T}_{\M} (f)^{(j)}(0)|)_{j\in\NN_0}$
%%	$$ (f^{(j)}(0)) R_j:=\mathcal{T}_{\M} (f)^{(j)}(0),$$
%	is equivalent to $\LL\M$.
%	Consequently, $\mathcal{T}_{\M} (f)$ is characteristic in the class $\mathcal{A}_{\{\LL\M\}} (S)$.
%\end{theorem}
%
%\demo{Proof}
%The first assertion is clear from Lemma~\ref{Scomplemma} and~\eqref{eq.DerivAt0T_Mf}. The second one stems from Theorem~\ref{optthm2}.
%%Since $(|f^{(j)}(0)|)_{j\in\NN_0}$ is equivalent to $\LL$ and by Lemma \ref{Scomplemma} we get that
%%$$\exists\;C,A\le 1\;\forall\;n\in\NN_0:\;\;\;|(f^{(j)}(0)) R_n |\geq  C A^n L_n M_n,$$
%%and the converse estimate follows analogously. Thus $(|(f^{(j)}(0))R_j|)_{j\in\NN_0}$ is equivalent to $\LL\M$.
%\qed\enddemo

\subsection{Construction of characteristic functions}\label{ssect.ConstrCharactFunct}
Given a sequence $\M\in\RR_{>0}^{\NN_0}$ and $\alpha>0$ we construct now, under suitable assumptions, characteristic functions in $\widetilde{\mathcal{A}}^u_{\{\M\}} (S_{\alpha})$.
% and in $\widetilde{\mathcal{A}}^u_{\{\M\}} (S_{\alpha})$.
For this we are using the basic functions from Subsection \ref{ssect.BasicFct} and the characteristic transform from Subsection \ref{ssect.CharactTransf}.%\vspace{6pt}
%
%
%\textbf{Uniform asymptotics.}
%

%{Let $\M\in\RR_{>0}^{\NN_0}$ and $\alpha>0$. We distinguish two situations: If $\alpha\le 2$, then take $\alpha':=\alpha$ and if $\alpha>2$, suppose that $\gamma(\M)+2>\alpha$ then take $\alpha':=\gamma(\M)+2$. Thus, in any case, we have that
	 %$$\overline{\G}^{2-\alpha'}\M:=(j^{(2-\alpha')j}M_j)_{j\in\NN_0}$$
	%is equivalent to a (lc) sequence $\LL$ depending on $\alpha'$, see Section \ref{ssect.IndexgammaM} and Theorem \ref{optthm1}.}

%\red{(Let $\M\in\RR_{>0}^{\NN_0}$ and $\alpha>0$. Assume that $\gamma(\M){+2}>\alpha$ and distinguish: If $\alpha\le 2$ then take $\alpha':=\alpha$ and if $\alpha>2$ then take $\alpha'$ such that $\gamma(\M){+2}>\alpha'>\alpha$. Thus, in any case, we have that
%	 $$\overline{\G}^{2-\alpha'}\M:=(j^{(2-\alpha')j}M_j)_{j\in\NN_0}$$
%	is equivalent to a {(lc)} sequence $\LL^{\alpha'}${(This notation is confused!?)} \red{belonging to the class \hyperlink{LCset}{$\mathcal{LC}$}}, see Section \ref{ssect.IndexgammaM} \red{(and Lemma)} {and Theorem \ref{optthm1}}.) }

%This requirement holds true if $M\in\hyperlink{LCset}{\mathcal{LC}}$ and $\alpha\le 2$; for $\alpha>2$ the previous condition can be replaced by the weaker assumption $\gamma(M)>\alpha$ since this requirement allows to switch to an equivalent sequence which belongs to the class \hyperlink{LCset}{$\mathcal{LC}$}.

\begin{theorem}\label{mainTthm1}
	Let $\M\in\RR_{>0}^{\NN_0}$ and $\alpha>0$.
	%
	\begin{enumerate}
		\item If $\alpha< 2$, we assume that  $\overline{\G}^{2-\alpha}\M:=(j^{(2-\alpha)j}M_j)_{j\in\NN_0}$
		is equivalent to an \hyperlink{lc}{$(\text{lc})$} sequence $\LL$. Then, $\mathcal{T}_{\LL} (\widetilde{E}_{\alpha}) $ is  characteristic in the class $\widetilde{\mathcal{A}}^u_{\{\M\}} (S_{\alpha})$.
		\item If $\alpha= 2$, we assume that  $\M$
		is equivalent to an \hyperlink{lc}{$(\text{lc})$} sequence $\LL$. Then, $\mathcal{T}_{\LL} (K_{RS}) $ is characteristic in the class $\widetilde{\mathcal{A}}^u_{\{\M\}} (S_{2})$.
		\item If $\alpha>2$, we assume that there exists $\a'>\a$ such that $\overline{\G}^{2-\alpha'}\M:=(j^{(2-\alpha')j}M_j)_{j\in\NN_0}$
		is equivalent to an \hyperlink{lc}{$(\text{lc})$} sequence $\LL$. Then, $\mathcal{T}_{\LL} (f_{\alpha,\alpha'}) $ is characteristic in the class $\widetilde{\mathcal{A}}^u_{\{\M\}} (S_{\alpha})$.
	\end{enumerate}
\end{theorem}
%
%
%
\demo{Proof}
This follows by Theorems~\ref{Ealphathm}, \ref{RStheorem}, \ref{fgalphaalphaprimethm},  \ref{CharacteristicTransformUniformAsymptotic} and \ref{CharacteristicTransformUniformAsymptotic1}, and  from the fact that $\overline{\G}^{\alpha-2}\LL$ in case 1, resp. $\LL$ in case 2, resp. $\overline{\G}^{\alpha'-2}\LL$ in case 3, is equivalent to $\M$.
\qed\enddemo
\begin{remark}%\label{rem.gammaindexcharactfunct}
	As it was carefully explained in~\cite[Remark 4.10]{ultraholomstab}), in order to guarantee that the hypotheses in the previous theorem are satisfied one can compute the lower Matuszewska index  $\gamma(\M)$, associated to the sequence of quotients $\m$ in the general theory of $O$-regular variation of sequences, and check whether it is greater than $\a-2$.
    % If this is the case, the very definition of this index implies that for {any $\beta$ such that $\gamma(\M)>\beta>\a-2$ the property $\left(P_{\beta}\right)$} (see Subsection~\ref{ssect.IndexgammaM}) is satisfied, and so there exists a suitable  (lc) sequence $\LL$ in the desired conditions.
\end{remark}


%(Detallar más???)

\begin{remark}\label{rem.signofcoefficients} For future use, observe that for any of the characteristic functions $f$ constructed in Theorem~\ref{mainTthm1}, if  $f\sim^u_{\{\M\}} \sum_{j=0}^\infty c_j z^j$, the coefficients are real and have alternating signs, in fact they satisfy $c_j=(-1)^j|c_j|$.
\end{remark}




%%%%%%%%%%%%%%%%%%%%%%%%%%%%%%%%%%%%%%%%%%%%%%%%%%%%%%%%%%%%%%%%%%%%%%%%%%%%%%
\section{Auxiliary formulas for point-wise multiplication and composition of asymptotic expansion}\label{sec.formula}

Following the approach in \cite{AubersonMennessier81}, we establish formulas describing the behavior of coefficients and remainders under point-wise multiplication and composition for functions, admitting a uniform asymptotic expansion at the point 0. 
%and defined on a sector region $S$ such that $0 \in \overline{S}$. 
Since we are dealing with different functions, in order to emphasize the dependence on the considered function we write for the appearing coefficients in the asymptotic expansions $c_j^f$, $c_j^g$ and so on. 

On the one hand, let $f$ and $g$ be functions defined in an arbitrary sector $S\subset \mathcal{R}$, and write $f(z)=\sum_{j=0}^{n-1}c^f_jz^j+r_f(z,n)$ and $g(z)=\sum_{j=0}^{n-1}c^g_jz^j+r_g(z,n)$. Then for the  point-wise multiplication $h:=f\cdot g$ we get $h(z)=\sum_{j=0}^{n-1}c^h_jz^j+r_h(z,n)$ with
\begin{equation}\label{eq.stablepointwisepropequ1}
c^h_j=\sum_{k=0}^jc^f_kc^g_{j-k},\hspace{15pt}r_h(z,n)=r_f(z,n)g(z)+\sum_{k=0}^{n-1}c^f_kz^kr_g(z,n-k);
\end{equation}
see also \cite[$(33)$]{AubersonMennessier81}.  

We also include, following~\cite[$(A.3)$]{AubersonMennessier81}, a version of the second equality in~\eqref{eq.stablepointwisepropequ1} for the case that $h=f^j=f\cdot f^{j-1}$ for a function $f$ such that $c_0^f=0$, and for any $n>j$. We write
\begin{align*}
(f(z))^j&=f(z)(f(z))^{j-1}=\left(\sum_{i=1}^{n-j}c^f_i z^i+r_f(z,n-j+1)\right)(f(z))^{j-1}\\
		&=\sum_{i=1}^{n-j}c^f_i z^i\left[\sum_{\ell=j-1}^{n-i-1}c^{f^{j-1}}_{\ell}z^{\ell}+r_{f^{j-1}}(z,n-i)\right]+r_f(z,n-j+1)(f(z))^{j-1}.
	\end{align*}
Since $f(z)=r_f(z,1)$, one deduces that
	\begin{equation}\label{eq.remaindern.fpowerj}
		 r_{f^j}(z,n)=\sum_{i=1}^{n-j}c^f_iz^ir_{f^{j-1}}(z,n-i)+r_f(z,n-j+1)\left(r_f(z,1)\right)^{j-1}.
	\end{equation}

On the other hand, let $f$ be a function defined in an arbitrary sector $S\subset \mathcal{R}$ with $c_0^f=0$, and $g$ be a function defined in a region $T\subset \CC$ containing the range of $f$. For $\varphi:=g\circ f$ we can write for all $z\in S$:
	$$\varphi(z)=c^g_0+\sum_{k=1}^{n-1}c^g_k(f(z))^k+r_g(f(z),n).$$
	Moreover, we have that
	\begin{equation*}%\label{equ70}
		(f(z))^k=\sum_{j=k}^{n-1}c^{f^k}_jz^j+r_{f^k}(z,n),
	\end{equation*}
	with
	\begin{equation}\label{equ71}
		c^{f^k}_j=\sum_{m_1+\dots+m_k=j}c^{f}_{m_1}\cdots c^{f}_{m_k},
	\end{equation}
	see \cite[$(70)$ \& $(71)$]{AubersonMennessier81}. Next we recall the crucial formulas \cite[$(73)$, $(74)$ \& $(75)$]{AubersonMennessier81}:
	\begin{equation*}%\label{equ73}
		\varphi(z)=\sum_{j=1}^{n-1}c^\varphi_jz^j+r_{\varphi}(z,n),
	\end{equation*}
	\begin{equation}\label{equ74}
		c^\varphi_0=c^g_0,\hspace{15pt}c^\varphi_n=\sum_{j=1}^{n}c^g_jc^{f^j}_n,\;\;\;n\in\NN,
	\end{equation}
	and
	\begin{equation}\label{eq.equ75.roumieu}
		r_\varphi(z,n)=\sum_{j=1}^{n-1}c^g_jr_{f^j}(z,n)+r_g(f(z),n).
	\end{equation}



\section{Stability under point-wise multiplication}\label{sec.alg}

We start with the following generalization of \cite[Thm. 3]{AubersonMennessier81}:


\begin{proposition}\label{prop.stablepointwiseprop}
Let $\M=(M_j)_j\in\RR_{>0}^{\NN_0}$ be a sequence with $M_0=1$, and let $S$ be a sector. If $\M$ satisfies \hyperlink{s-alg}{$(\on{alg})$}, then the class $\mathcal{\widetilde{A}}^{u}_{[\M]}(S)$ is closed under point-wise multiplication of functions; i.e. $\mathcal{\widetilde{A}}^{u}_{[M]}(S)$ is an algebra. 
\end{proposition}

\begin{proof} Since $\M$ satisfies \hyperlink{s-alg}{$(\on{alg})$}, we have $M_jM_k\le C^{j+k}M_{j+k}$ for some $C\geq 1$ and for all $j,k\in\NN_0$.
Let $f,g\in\mathcal{\widetilde{A}}^{u}_{[\M]}(S)$.  Now, by the definition of the classes, we have $f\in\mathcal{\widetilde{A}}^{u}_{\M,d_1}(S)$ and $g\in\mathcal{\widetilde{A}}^{u}_{\M,d_2}(S)$ for some (resp. all) $d_i>0$, and so there exist $A_1,A_2>0$ such that
$$
|r_f(z,n)|\le  A_1 d_1^n M_n|z|^n,\quad 
|r_g(z,n)|\le  A_2 d_2^n M_n|z|^n,\quad z\in S.
$$ 
Thus by \eqref{eq.stablepointwisepropequ1} and \eqref{cjestimation}, we estimate as follows for all $n\in\NN_0$ and $z\in S$:

\begin{align*}
|r_h(z,n)|&\le |r_f(z,n)||g(z)|+\sum_{k=0}^{n-1}|c^f_k||z|^k|r_g(z,n-k)|\\&
\le A_1 d_1^n M_n|z|^n A_2 + \sum_{k=0}^{n-1} A_1 d_1^k M_k|z|^k A_2 d_2^{n-k}M_{n-k}|z|^{n-k}
\\&
\le A_1A_2d_1^nM_n|z|^n + A_1A_2\sum_{k=0}^{n-1} d_1^k d_2^{n-k}C^nM_n|z|^n
\\&
\le A_1A_2d_1^nM_n|z|^n+A_1A_2C^nM_n|z|^n\sum_{k=0}^{n}\binom{n}{k}d_1^kd_2^{n-k}
\\&
\le A_1A_2d_1^nM_n|z|^n+ A_1A_2C^nM_n|z|^n(d_1+d_2)^n
\\&
\le 2A_1A_2(C(d_1+d_2))^nM_n|z|^n.
\end{align*}
This finishes the Roumieu part. Concerning the Beurling class $\mathcal{\widetilde{A}}^{u}_{(\M)}$ we use the same estimate, and notice that $C(d_1+d_2)\rightarrow 0$ when $d_i\rightarrow 0$ since $C$ is only depending on $\M$ via \hyperlink{s-alg}{$(\on{alg})$} but not on $d_i$, $i=1,2$.
\end{proof}

Conversely, if we can ensure the existence of very well-suited characteristic functions in the Roumieu class, closure under point-wise multiplication implies condition \hyperlink{s-alg}{$(\on{alg})$}. The following result gathers different implications.

\begin{theorem}\label{thm.stablepointwiseprop}
Let $\M=(M_j)_j\in\RR_{>0}^{\NN_0}$ be a sequence with $M_0=1$ and $\alpha>0$. Consider the following statements:
\begin{itemize}
    \item[(i)] In case $\alpha\leq 2$, $\overline{\G}^{2-\alpha}\M$ is equivalent to an \hyperlink{lc}{$(\text{lc})$} sequence; if $\alpha>2$, there exists $\a'>\a$ such that $\overline{\G}^{2-\alpha'}\M$ is equivalent to an \hyperlink{lc}{$(\text{lc})$} sequence.
\item[(ii)] $\M$ satisfies \hyperlink{s-alg}{$(\on{alg})$}.
\item[(iii)] $\mathcal{\widetilde{A}}^{u}_{\{\M\}}(S_{\a})$ is an algebra.
\item[(iv)] There exists a characteristic function $f\in \mathcal{\widetilde{A}}^{u}_{\{\M\}} (S_\alpha)$, with  $f\sim^u_{\{\M\}} \sum_{j=0}^\infty c^f_j z^j$, such that $(|c^f_j|)_{j\in\NN_0}$ is equivalent to $\M$ and, moreover,  $c^f_j=(-1)^j|c^f_j|$ for all $j\in\NN_0$.
\end{itemize} 
Then, one has $(i)\Rightarrow(ii)\Rightarrow(iii)$, $(i)\Rightarrow(iv)$ and $(iii)+(iv)\Rightarrow(ii)$. 
\end{theorem}
\begin{proof}
$(i)\Rightarrow(ii)$ This a consequence of the following general statement, with a straightforward proof: Assume that we have sequences $\mathbf{M}$, $\mathbf{N}$ (with $M_0=1=N_0$) such that $\mathbf{M}\cdot\mathbf{N}\approx\mathbf{L}$ and $\mathbf{L}$ satisfies \hyperlink{s-alg}{$(\on{alg})$} (and $L_0=1$). If $\mathbf{N}$ satisfies the condition of moderate growth, namely there exists $A>0$ such that $N_{j+k}\le A^{j+k}N_jN_k$ for all $j,k\in\NN_0$, then $\mathbf{M}$ satisfies \hyperlink{s-alg}{$(\on{alg})$} as well. In our case, it suffices to observe that \hyperlink{lc}{$(\text{lc})$} implies \hyperlink{s-alg}{$(\on{alg})$} and that the Gevrey-like sequences $\overline{\G}^{\beta}$ of every order $\beta\in\RR$ satisfy the moderate growth condition.

% For $\mathbf{L}$ we get the estimate $L_jL_k\le C^{j+k}L_{j+k}$ for some $C\ge 1$ and all $j,k\in\NN_0$ and the equivalence yields for a constant $A\ge 1$ that $M_jN_jM_kN_k\le A^jL_jA^kL_k\le(AC)^{j+k}L_{j+k}\le(A^2C)^{j+k}M_{j+k}N_{j+k}$ for all $j,k\in\NN_0$. Now, since $\mathbf{N}$ has (mg) we obtain $N_{j+k}\le B^{j+k}N_jN_k$ for some $B\ge 1$ and all $j,k\in\NN_0$ and, summarizing,
% $$M_jM_k\le(A^2C)^{j+k}M_{j+k}\frac{N_{j+k}}{N_jN_k}\le(A^2BC)^{j+k}M_{j+k},$$
% which is (alg) for $\mathbf{M}$ with the constant $A^2BC$.

$(ii)\Rightarrow(iii)$ This is  Proposition~\ref{prop.stablepointwiseprop}.

$(i)\Rightarrow(iv)$ See Theorem~\ref{mainTthm1} and Remark~\ref{rem.signofcoefficients}.

$(iii)+(iv)\Rightarrow(ii)$ Let $f\in \mathcal{\widetilde{A}}^{u}_{\{\M\}} (S_\alpha)$ be as in $(iv)$. Since $\mathcal{\widetilde{A}}^{u}_{\{\M\}} (S_\alpha)$ is an algebra we have that $h= f\cdot f \in \mathcal{\widetilde{A}}^{u}_{\{\M\}} (S_\alpha)$.

By \eqref{cjestimation}, there exist $A_1>0$ and $d_1>0$ such that $|c^h_n|\le A_1d_1^nM_n$ for all $n\in\NN_0$. Using the auxiliary formula for the coefficients of a product \eqref{eq.stablepointwisepropequ1}, we have $c^h_n = \sum_{k=0}^n c^f_k c^f_{n-k}$. So, for any $j\in\{0,1,2,\dots,n\}$ we see that
$$|c^h_n|=|\sum_{k=0}^n c^f_k \cdot c^f_{n-k}| = |\sum_{k=0}^n  (-1)^k|c^f_k| \cdot (-1)^{n-k} |c^f_{n-k}|| = \sum_{k=0}^n |c^f_k| \cdot |c^f_{n-k}|  \ge |c^f_j| \cdot |c^f_{n-j}|.$$
Since $(|c^f_j|)_{j\in\NN_0}$ is equivalent to $\M$, there exist $A_2,d_2>0$ such that $|c^f_j|\geq A_2 d_2^j M_j$ and for all $j,k\in\NN_0$, by taking $n=j+k$, we deduce that 
 $$A_1d_1^{j+k}M_{j+k} \geq |c^h_{j+k}| \ge |c^f_j| \cdot |c^f_{k}|\ge A_2^2 d_2^{j+k} M_j M_k.$$
 Consequently, $\M$ satisfies \hyperlink{s-alg}{$(\on{alg})$}.
\end{proof}
 


% \section{Stability properties of uniform asymptotic expansion\\ classes defined by weight matrices}\label{provsect.StabilityPropert}


% \begin{proposition}\label{stablepointwiseprop}
% Let $\mathcal{M}=\{\M^{(p)}: p>0\}$ be a weight matrix satisfying \hyperlink{R-alg}{$(\mathcal{M}_{\{\on{alg}\}})$}. Then the class $\mathcal{\widetilde{A}}^{u}_{\{\mathcal{M}\}}$ is stable under point-wise multiplication of functions; i.e. $\mathcal{\widetilde{A}}^{u}_{\{\mathcal{M}\}}$ is an algebra. If $\mathcal{M}$ satisfies \hyperlink{B-alg}{$(\mathcal{M}_{(\on{alg})})$}, then $\mathcal{\widetilde{A}}^{u}_{(\mathcal{M})}$ is an algebra.
% \end{proposition}

% \begin{proof}
% Let $f,g\in\mathcal{\widetilde{A}}^{u}_{\{\mathcal{M}\}}(S)$, $S$ an arbitrary sector, and write $f(z)=\sum_{j=0}^{n-1}c^f_jz^j+r_f(z,n)$ and $g(z)=\sum_{j=0}^{n-1}c^g_jz^j+r_g(z,n)$. Then for $h:=f\cdot g$ we get $h(z)=\sum_{j=0}^{n-1}c^h_jz^j+r_h(z,n)$ with
% \begin{equation}\label{stablepointwisepropequ1}
% c^h_j=\sum_{k=0}^jc^f_kc^g_{j-k},\hspace{15pt}r_h(z,n)=\sum_{k=0}^{n-1}c^f_kz^kr_g(z,n-k)+r_f(z,n)g(z);
% \end{equation}
% see also \cite[$(33)$]{AubersonMennessier81}. Now, by definition of the classes we have $f\in\mathcal{\widetilde{A}}^{u}_{\M^{(p_1)},h_1}(S)$ and $g\in\mathcal{\widetilde{A}}^{u}_{\M^{(p_2)},h_2}(S)$ for some $p_i,h_i>0$. We put $\overline{p}:=\max\{p_1,p_2\}$ and apply  \hyperlink{R-alg}{$(\mathcal{M}_{\{\on{alg}\}})$} to this index. Thus, $M^{(\overline{p})}_jM^{(\overline{p})}_k\le C^{j+k}M^{(q)}_{j+k}$ for some index $q$, some $C>0$ and all $j,k\in\NN_0$. W.l.o.g. we can take $q\ge\overline{p}$ and $C\ge 1$. Thus, by taking into account this fact, \eqref{stablepointwisepropequ1} and \eqref{cjestimation}, we estimate as follows for all $n\in\NN_0$ and $z\in S$:

% \begin{align*}
% |r_h(z,n)|&\le\sum_{k=0}^{n-1}|c^f_k||z|^k|r_g(z,n-k)|+|r_f(z,n)||g(z)|
% \\&
% \le\sum_{k=0}^{n-1}A_1h_1^kM^{(p_1)}_k|z|^kA_2h_2^{n-k}M^{(p_2)}_{n-k}|z|^{n-k}+A_1h_1^nM^{(p_1)}_n|z|^nA_2
% \\&
% \le A_1A_2\sum_{k=0}^{n-1}h_1^kh_2^{n-k}M^{(\overline{p})}_kM^{(\overline{p})}_{n-k}|z|^n+A_1A_2h_1^nM^{(\overline{p})}_n|z|^n
% \\&
% \le A_1A_2\sum_{k=0}^{n-1}h_1^kh_2^{n-k}C^nM^{(q)}_n|z|^n+A_1A_2h_1^nM^{(q)}_n|z|^n
% \\&
% \le A_1A_2C^nM^{(q)}_n|z|^n\sum_{k=0}^{n}\binom{n}{k}h_1^kh_2^{n-k}+A_1A_2h_1^nM^{(q)}_n|z|^n
% \\&
% \le A_1A_2C^nM^{(q)}_n|z|^n(h_1+h_2)^n+A_1A_2h_1^nM^{(q)}_n|z|^n
% \\&
% \le 2A_1A_2(C(h_1+h_2))^nM^{(q)}_n|z|^n.
% \end{align*}
% This finishes the Roumieu part. Concerning the Beurling class $\mathcal{\widetilde{A}}^{u}_{(\mathcal{M})}$ we use the same estimate, apply \hyperlink{B-alg}{$(\mathcal{M}_{(\on{alg})})$} and notice that $C(h_1+h_2)\rightarrow 0$ when $h_i\rightarrow 0$ since $C$ is only depending on the chosen index $q$ (via \hyperlink{B-alg}{$(\mathcal{M}_{(\on{alg})})$}) but not on $h_i$.
% \end{proof}

% \begin{remark}\label{generalindexremark}
% Note that the proof of Proposition \ref{stablepointwiseprop} shows by iteration that when $f$ requires the choice $\alpha_1$, then for $f^n$ it suffices to take $\alpha_n$ and the sequence of indices $(\alpha_n)_n$ is given iteratively via \hyperlink{R-alg}{$(\mathcal{M}_{\{\on{alg}\}})$}; i.e. $\alpha_{n-1}=\alpha$ and $\alpha_{n}=\beta$ in this condition. In particular, $n\mapsto\alpha_n$ is non-decreasing.
% \end{remark}




\section{Stability under composition}\label{sect.StabilityCompos}

Let us first define a suitable concept of stability under composition for classes associated with the sequence $\M$.

\begin{definition}\label{def.UniformAsympComplexPlane}
Let $V$ be an open subset in the complex plane such that $0\in\overline{V}$, $f\colon V\to\CC$ be a holomorphic function on $V$, $\M=(M_j)_j\in\RR_{>0}^{\NN_0}$ be a sequence with $M_0=1$, and $h>0$. We say $f\in\widetilde{\mathcal{A}}_{\M,h}^{u}(V)$ if there exists a formal complex power series $\sum_{k=0}^{\infty}c_kz^k$  such that
\[
%\left\|f\right\|_{\M,h,\overset{\sim}{u}}:=
\sup_{z\in V\setminus\{0\},\,j\in\NN_{0}}\frac{|f(z)-\sum_{k=0}^{j-1}c_kz^k|}{h^{j}M_{j}|z|^j}<+\infty.
\]   
We set 
$$\widetilde{\mathcal{A}}_{\{\M\}}^{u}(V):=\bigcup_{h>0}\widetilde{\mathcal{A}}_{\M,h}^{u}(V),\qquad
\widetilde{\mathcal{A}}_{(\M)}^{u}(V):=\bigcap_{h>0}\widetilde{\mathcal{A}}_{\M,h}^{u}(V).$$
%\end{equation}
\end{definition}

Observe that, in case $V$ is a sector $S$ with vertex at 0 in the complex plane, this definition agrees with that given for $\widetilde{\mathcal{A}}_{[\M]}^{u}(S)$
in Section~\ref{sect.AsymptoticexpClass}. 

\begin{definition}\label{def.StabilityCompos}
Given a sequence $\M$ and a sector $S$, we say the class $\widetilde{\mathcal{A}}_{[\M]}^{u}(S)$ is stable under composition if for every $f\in\widetilde{\mathcal{A}}_{[\M]}^{u}(S)$ with $c^f_0=0$ (with the notation in Section~\ref{sec.formula}), for every subsector $T$ of $S$ and for every $g\in\widetilde{\mathcal{A}}_{[\M]}^{u}(f(T))$ one has that $g\circ (f|_T)\in\widetilde{\mathcal{A}}_{[\M]}^{u}(T)$.
\end{definition}

\begin{remark}
In the previous definition, the condition $c_0^f=0$ guarantees, thanks to the expression~\eqref{eq.coeff.asymp.expan.lim.remainder} for $n=0$, that $0\in\overline{f(T)}$ for every subsector $T$ of $S$, and so it makes sense to consider functions $g$ in $\widetilde{\mathcal{A}}_{[\M]}^{u}(f(T))$. 
\end{remark}

The next auxiliary result studies the growth of the remainders for the powers $f^j$, $j\in\NN$, for a function $f\in\widetilde{\mathcal{A}}_{\M,h}^{u}(S)$ with $c_0^f=0$. We note first that, if $n\in\NN_0$ and $n\le j$, the fact that $c_0^f=0$ implies that $r_{f^j}(z,n)=f^j(z)$ for $z\in S$. Since $f(z)=r_f(z,0)=r_f(z,1)$, we know that $|f(z)|\le \min\{A,AhM_1|z|\}$ for every $z\in S$ and some $A>0$, and so
\begin{equation}\label{eq.bound.remainders.fpowerj.n_at_most_j}
    |r_{f^j}(z,n)|\le A^j(\min\{1,hM_1|z|\})^j,\quad z\in S.
\end{equation}
We treat now the case $n>j\ge 1$.

\begin{lemma}\label{lemma.bounds.powers}
Let $\M=(M_j)_j\in\RR_{>0}^{\NN_0}$ be a sequence with $M_0=1$, $S$ be a sector and $f\in\widetilde{\mathcal{A}}_{\M,h}^{u}(S)$ with $c^f_0=0$, so that there exists $A>0$ such that
\begin{equation}\label{eq.bounds.fpower1}
    |r_f(z,n)|\le A h^n M_n |z|^n,\quad z\in S,\ n\in\NN_0.
\end{equation}
Then, for every $j,n\in\NN$ with $n>j$ one has
\begin{equation}\label{eq.bounds.fpowerj}
    |r_{f^j}(z,n)|\le A^j h^n \left(\sum_{\substack{m_1+\dots+m_j=n\\m_1\ge 1,\dots,m_j\ge 1}}M_{m_1}\cdot\ldots \cdot M_{m_j}\right) |z|^n,\quad z\in S.
\end{equation}
\end{lemma}
\begin{remark}
Observe that the estimates~\eqref{eq.bounds.fpowerj} are also valid when $n=j$, since they take the form
$$
|r_{f^j}(z,j)|\le A^jh^jM_1^j|z|^j,\quad z\in S,
$$
which was already obtained in~\eqref{eq.bound.remainders.fpowerj.n_at_most_j}.
\end{remark}
\begin{proof}
We proceed by induction on $j$. For $j=1$ it is clear that the estimates in~\eqref{eq.bounds.fpowerj} reduce to those in~\eqref{eq.bounds.fpower1} particularized for $n>1$. Suppose now that the estimates~\eqref{eq.bounds.fpowerj} hold true for some $j-1\ge 1$ and every $n>j-1$, and we treat the $j$-th power. 
Observe that, whenever $1\leq i \leq n-j$, one has $n-i\geq n-(n-j)>j-1$, and so we can apply the induction hypothesis in order to estimate $r_{f^{j-1}}(z,n-i)$ in the formula~\eqref{eq.remaindern.fpowerj}. Moreover, we can use~\eqref{cjestimation} and write
\begin{align*}
|r_{f^j}(z,n)|&\leq \sum_{i=1}^{n-j}Ah^iM_i|z|^iA^{j-1}h^{n-i}\left(\sum_{\substack{m_1+\dots+m_{j-1}=n-i\\ m_1\ge 1,\dots,m_{j-1}\ge 1}}M_{m_1}\cdot\ldots \cdot M_{m_{j-1}}\right)|z|^{n-i}\\
&+Ah^{n-j+1}M_{n-j+1}|z|^{n-j+1}(AhM_1|z|)^{j-1}\\
&=A^j h^n \sum_{i=1}^{n-j+1}M_i\left(\sum_{\substack{m_1+\dots+m_{j-1}=n-i\\m_1\ge 1,\dots,m_{j-1}\ge 1}}M_{m_1}\cdot\ldots \cdot M_{m_{j-1}}\right) |z|^n\\
&=A^j h^n \left(\sum_{\substack{m_1+\dots+m_j=n\\m_1\ge 1,\dots,m_j\ge 1}}M_{m_1}\cdot\ldots \cdot M_{m_j}\right) |z|^n,
	\end{align*}
as desired.
\end{proof}

We are ready for the following generalization of \cite[Thm. 6]{AubersonMennessier81}.


\begin{proposition}\label{prop.stablecomposition}
Let $\M=(M_j)_j\in\RR_{>0}^{\NN_0}$ be a sequence with $M_0=1$, and let $S$ be a sector. If $\M$ satisfies \hyperlink{s-FdB}{$(\on{FdB})$}, then the class $\mathcal{\widetilde{A}}^{u}_{[\M]}(S)$ is closed under composition. 
\end{proposition}
\begin{proof}
Let $f\in\mathcal{\widetilde{A}}^{u}_{[\M]}(S)$ with $c_0^f=0$, $T$ be a subsector of $S$, and  $g\in\mathcal{\widetilde{A}}^{u}_{[\M]}(f(T))$. We have $f|_T\in\mathcal{\widetilde{A}}^{u}_{\M,h_1}(T)$ and $g\in\mathcal{\widetilde{A}}^{u}_{\M,h_2}(f(T))$ for some (resp. all) $h_i>0$, $i=1,2$, so that there exist $A_1=A_1(h_1)>0$ and $A_2=A_2(h_2)>0$ such that
\begin{equation*}%\label{eq.bounds.f_h1}
    |r_f(z,n)|\le A_1h_1^n M_n |z|^n,\quad z\in T,\ n\in\NN_0,
\end{equation*}
and
\begin{equation*}%\label{eq.bounds.g_h2}
    |r_g(f(z),n)|\le A_2 h_2^n M_n |f(z)|^n,\quad z\in T,\ n\in\NN_0.
\end{equation*}
Put $\varphi=g\circ f|_T$; according to the formula~\eqref{eq.equ75.roumieu}, and subsequently using Lemma~\ref{lemma.bounds.powers} and the fact that $f(z)=r_f(z,1)$, for every $n\in\NN_0$ and $z\in T$ we have
\begin{align*}
|r_{\varphi}(z,n)|&\le\sum_{j=1}^{n-1}|c^g_j||r_{f^j}(z,n)|+|r_g(f(z),n)|\\
&\le \sum_{j=1}^{n-1} A_2 h_2^j M_j \cdot A_1^j h_1^n 
\left(\sum_{\substack{m_1+\dots+m_j=n\\m_1\ge 1,\dots,m_j\ge 1}}M_{m_1}\cdot\ldots \cdot M_{m_j}\right) |z|^n
\\
&\quad + A_2 h_2^n M_n \left( A_1 h_1 M_1 |z| \right)^n\\
&=\sum_{j=1}^{n} A_2 h_2^j M_j \cdot A_1^j h_1^n 
\left(\sum_{\substack{m_1+\dots+m_j=n\\m_1\ge 1,\dots,m_j\ge 1}}M_{m_1}\cdot\ldots \cdot M_{m_j}\right) |z|^n.
\end{align*}
By assumption, $\M$ satisfies \hyperlink{s-FdB}{$(\on{FdB})$}, so there exist $C,B>0$ such that, for any $m_1,\dots,m_j\in\NN$ with $m_1+\dots+m_j=n$, one has 
$$
M_jM_{m_1}\cdot\ldots \cdot M_{m_j}\le CB^nM_n.
$$
Since the number of compositions of $n$ as the sum of $j$ positive integers is  $\binom{n-1}{j-1}$, we can write
\begin{align*}
|r_{\varphi}(z,n)|
&\le C A_2 (Bh_1)^n M_n |z|^n 
\sum_{j=1}^{n} \binom{n-1}{j-1} (A_1 h_2)^j\\
&\le C A_2 \left( Bh_1(1 + A_1 h_2) \right)^n M_n |z|^n .
\end{align*}
So, we deduce that $\varphi\in\mathcal{\widetilde{A}}^{u}_{\M,h_3}(T)$ for $h_3=Bh_1 (1 + A_1 h_2)$, and we are done in the Roumieu case. In the Beurling case, given an arbitrary $h_3>0$ one can take for $f$ the value $h_1>0$ such that $Bh_1(1+h_1)=h_3$, which will provide a value $A_1>0$ associated to $h_1$ for $f$; then, one can choose for $g$ the value $h_2=h_1/A_1>0$ and we conclude.
\end{proof}

Conversely and as before, whenever we have suitable characteristic functions in the Roumieu class, closure under composition implies condition \hyperlink{s-FdB}{$(\on{FdB})$}. 

\begin{theorem}\label{thm.stablecomposition}
Let $\M=(M_j)_j\in\RR_{>0}^{\NN_0}$ be a sequence with $M_0=1$ and $\alpha>0$. If $\alpha\leq 2$ we assume that $\overline{\G}^{2-\alpha}\M$ is equivalent to an \hyperlink{lc}{$(\text{lc})$} sequence, and if $\alpha>2$ we assume that there exists $\a'>\a$ such that $\overline{\G}^{2-\alpha'}\M$ is equivalent to an \hyperlink{lc}{$(\text{lc})$} sequence. Then the class $\mathcal{\widetilde{A}}^{u}_{\{\M\}} (S_\alpha)$ is closed under composition if and only if $\M$ satisfies  \hyperlink{s-FdB}{$(\on{FdB})$}. 
\end{theorem}
 \begin{proof} By Proposition~\ref{prop.stablecomposition}, we only need to show necessity. 
 
By Theorem~\ref{mainTthm1}, there exists a characteristic function $f\in \mathcal{\widetilde{A}}^{u}_{\{\M\}} (S_\alpha)$ in the strong sense, i.e., $f\sim^u_{\{\M\}} \sum_{j=0}^\infty c^f_j z^j$ and $(|c^f_j|)_{j\in\NN_0}$ is equivalent to $\M$. Moreover, by Remark~\ref{rem.signofcoefficients} we know that $c_0^f>0$ and $c_1^f<0$. Let us consider the function $f_0$ defined in $S_\a$ by $f_0(z)=c_0^f-f(z)$. It is clear that $f_0\in \mathcal{\widetilde{A}}^{u}_{\{\M\}} (S_\alpha)$ and $f_0\sim^u_{\{\M\}} -\sum_{j=1}^\infty c^f_j z^j$ (in particular, $c_0^{f_0}=0$). 

In case $\alpha\le 2$, choose real numbers $\b$ and $\varepsilon$ such that $0<\varepsilon<\frac{\pi}{2}(\alpha-\b)<\frac{\pi}{2}$. Consider $\varepsilon_1:=-c_1^f\sin(\varepsilon)<|c_1^f|$.
Due to~\eqref{eq.limit.derivatives.coeff}, we have that
$$
\lim_{\substack{z\to 0\\z\in S_{\b}}}f_0'(z)=-c_1^f,
$$
and so there exists $r>0$ such that for every $z\in S_{\b,r}=\{w\in S_{\b}\colon |w|<r\}$ one has $|f_0'(z)+c_1^f|<\varepsilon_1$. Take $z\in S_{\b,r}$, then the segment $(0,z]$ is contained in $S_{\b,r}$ and so,
$$
|f_0(z)+c_1^fz|=\left|\int_{[0,z]}(f_0'(w)+c_1^f)\,dw\right|\le \varepsilon_1|z|.
$$
For such $z$ we deduce that
$$
|\arg(f_0(z))-\arg(z)|=|\arg(f_0(z))-\arg(-c_1^fz)|\le \arcsin\left(\frac{\varepsilon_1|z|}{-c_1^f|z|}\right)=\varepsilon.
$$
The choices of $\b$ and $\varepsilon$ guarantee that $|\arg(f_0(z))|<\frac{\a\pi}{2}$, and so $f_0(S_{\b,r})\subset S_{\a}$. Since $\mathcal{\widetilde{A}}^{u}_{\{\M\}} (S_\alpha)$ is closed under composition, we know that $\varphi:= f\circ f_0 \in \mathcal{\widetilde{A}}^{u}_{\{\M\}} (S_{\b,r})$ (where, for simplicity, we have written $f_0$ instead of $f_0|_{S_{\b,r}}$, and $f$ instead of $f|_{f_0(S_{\b,r})}$).

In particular, if $\varphi\sim^u_{\{\M\}}\sum_{n=0}^\infty c_n^{\varphi}z^n$, we know that there exist $A_1>0$ and $h_1>0$ such that 
\begin{equation}\label{eq.estimates.cn.composition}
|c^{\varphi}_n|\le A_1h_1^nM_n, \quad n\in\NN_0.     
\end{equation}
At the same time, the formulas~\eqref{equ71} and~\eqref{equ74} allow us to write $c^\varphi_0=c^f_0$ and, since $c_0^{f_0}=0$ and using again Remark~\ref{rem.signofcoefficients},
\begin{align*}
c^\varphi_n&=\sum_{j=1}^{n}c^f_jc^{f_0^j}_n=
\sum_{j=1}^{n}c^f_j\sum_{m_1+\dots+m_j=n}c^{f_0}_{m_1}\cdots c^{f_0}_{m_j}\\
&=\sum_{j=1}^{n}(-1)^j|c^f_j|\sum_{\substack{m_1+\dots+m_j=n\\m_1\ge 1,\dots,m_j\ge 1}}(-(-1)^{m_1}|c^{f}_{m_1}|)\cdots (-(-1)^{m_j}|c^{f}_{m_j}|)\\
&=(-1)^n\sum_{j=1}^{n}|c^f_j|\sum_{\substack{m_1+\dots+m_j=n\\m_1\ge 1,\dots,m_j\ge 1}}|c^{f}_{m_1}|\cdots |c^{f}_{m_j}|,\;\;\;n\in\NN. 
\end{align*}
As $(|c^f_j|)_{j\in\NN_0}$ is equivalent to $\M$, there exist $A_2,B>0$ such that $|c^f_j|\geq A_2 B^j M_j$ for all $j\in\NN_0$. Then,
\begin{align*}
|c^\varphi_n|&=\sum_{j=1}^{n}|c^f_j|\sum_{\substack{m_1+\dots+m_j=n\\m_1\ge 1,\dots,m_j\ge 1}}|c^{f}_{m_1}|\cdots |c^{f}_{m_j}|\\
&\ge \sum_{j=1}^{n}A_2 B^j M_j\sum_{\substack{m_1+\dots+m_j=n\\m_1\ge 1,\dots,m_j\ge 1}}A_2 B^{m_1}M_{m_1}\cdots A_2 B^{m_j} M_{m_j}\\
&=A_2 B^n\sum_{j=1}^{n} (A_2B)^jM_j\sum_{\substack{m_1+\dots+m_j=n\\m_1\ge 1,\dots,m_j\ge 1}}M_{m_1}\cdots M_{m_j}.
\end{align*}
In case $A_2B\ge 1$, we deduce that, for every $m_1,\dots,m_j\in\NN$ with $\sum_{i=1}^jm_i=n$, we have $|c_n^{\varphi}|\ge A_2B^nM_jM_{m_1}\dots M_{m_j}$, and we obtain the condition \hyperlink{s-FdB}{$(\on{FdB})$} when comparing with the inequalities~\eqref{eq.estimates.cn.composition}. If $A_2B<1$, we get
$|c_n^{\varphi}|\ge A_2(A_2B^2)^nM_jM_{m_1}\dots M_{m_j}$, and the conclusion follows similarly.

Finally, the case $\a>2$ is treated by considering the restriction of a charasteristic function $f$ in the class $\mathcal{\widetilde{A}}^{u}_{\{\M\}} (S_\alpha)$ to the sector $S_2$, which can be identified with $\CC\setminus(-\infty,0]$, and proceeding as before to obtain a proper subsector $S_{\b,r}$, with $\b\in(1,2)$ and $r>0$, so that the image of $S_{\b,r}$ under the function $c_0^f-f$ is contained in $S_2$. This allows us to reason as before with the function $f\circ(c_0^f-f)$ (after suitable restrictions), which belongs to the corresponding class by hypothesis, and we can deduce the  condition \hyperlink{s-FdB}{$(\on{FdB})$} again.
\end{proof}



%%  esta llave cierra el cambio a color morado

% \section{Special cases}\label{specialcasesection}

% \subsection{The weight sequence case}\label{sect.StabilityWeightSequenceCase}
% For a sequence $\M\in\RR_{>0}^{\NN_{0}}$, we study when the class $\widetilde{\mathcal{A}}^{u}_{[\M]}(S)$ is closed under composition. We are able to conclude by applying Theorem~\ref{Theo.CompStability} and Theorem~\ref{Theo.CompStability.Beur} to the constant weight matrix $\mathcal{M}=\{\M^{(p)}=\M: p>0\}$ and hence get the following result:

% \begin{corollary}\label{Coro.CompStability.seq}
% 	Let $\M$ be a sequence and $S$ a sector. If $\M$ has \hypertarget{salg}{$(\on{salg})$}, \hypertarget{pri}{$(\on{pri})$} and $(p!)_p\subset \M$, then the class
% 	$\widetilde{\mathcal{A}}^{u}_{[\M]}(S)$ is closed under composition; i.e. for all $f\in\widetilde{\mathcal{A}}^{u}_{[\M]}(S)$ with $f(0)=0$ and for all $g\in\widetilde{\mathcal{A}}^{u}_{[\M]}(T)$, where $T\subset \mathcal{R}$ is a sector containing the range of $f$, we have $g\circ f\in\widetilde{\mathcal{A}}^{u}_{[\M]}(S)$.
% \end{corollary}

% Consequently, in view of Remark \ref{newcondrem} it follows that Corollary \ref{Coro.CompStability.seq} applies, in particular, to the following important and known examples of sequences:

% \begin{itemize}
% \item[$(*)$] $\mathbf{M}=\mathbf{G}^a$ for any $a\ge 1$ (Gevrey sequences with index $a\ge 1$);

% \item[$(*)$] $\mathbf{M}=(q^{p^2})_p$ with $q>1$ arbitrary ($q$-Gevrey sequences).
% \end{itemize}




% \subsection{The weight function case}\label{sect.StabilityWeightFUnctCase}

% The weight function case is now shown by involving the associated weight matrix $\mathcal{M}_{\omega}$. If $\omega\in\mathcal{W}$ is given, then recall that $\mathcal{M}_{\omega}$ always has \hyperlink{Mlc}{$(\mathcal{M}_{\on{lc}})$} and, in addition, it turns out that condition $(\omega_2)$ for $\omega$ implies \hyperlink{holom}{$(\mathcal{M}_{\mathcal{H}})$} for $\mathcal{M}_{\omega}$ (see subsection~\ref{ssect.WeightMatricesFromFunctions}).

%%%%%%%%%%%%%%%%%%%%%%%%%%%%%%%%%%%%%%%%%%%%%%%%
%
%    ESTO YA ESTABA COMENTADO
%
%We start proving, for the reader's convenience, how the condition $(\omega_2)$ for $\omega$ is sufficient in order to prove \hyperlink{sum}{$(\mathcal{M}_{\on{sum}})$} for the weight matrix associated to the weight function $\omega$.
%
%\begin{lemma}\label{sum-omega2}
%	Let $\omega\in\hyperlink{omset0}{\mathcal{W}_0}$ be given with associated weight matrix $\mathcal{M}_{\omega}:=\{\mathbb{W}^{(\ell)}: \ell>0\}$. If $\omega$ has $(\omega_2)$, then  $\mathcal{M}_{\omega}$ has \hyperlink{sum}{$(\mathcal{M}_{\on{sum}})$}.
%\end{lemma}
%\begin{proof}
%	If $\omega$ satisfies the condition $(\omega_2)$, then $\mathcal{M}_{\omega}$ satisfies \hyperlink{holom}{$(\mathcal{M}_{\mathcal{H}})$}. Moreover, thanks to the fact that $\mathcal{M}_{\omega}$ has \hyperlink{Mlc}{$(\mathcal{M}_{\on{lc}})$}, we deduce from Corollary~\ref{Coro.equivliminf} that $\liminf_{p\rightarrow\infty}w^{(\ell)}_p/p>0$, for all $\ell>0$. Now, by arguing for each row, we deduce from \hyperlink{Mlc}{$(\mathcal{M}_{\on{lc}})$}, the previous inferior limit condition and Theorem~\ref{aubsersonmennessiercondition}, that  $\mathcal{M}_{\omega}$ has \hyperlink{sum}{$(\mathcal{M}_{\on{sum}})$}.
%\end{proof}
%%%%%%%%%%%%%%%%%%%%%%%%%%%%%%%%%%%%%%%%%%%%%%%%%%

% We can provide now a statement about stability properties for classes associated to a weight function.

% \begin{theorem}\label{stability-weightfunction}
% 	Let $\omega\in\hyperlink{omset1}{\mathcal{W}}$ be given with associated weight matrix $\mathcal{M}_{\omega}:=\{\W^{(\ell)}: \ell>0\}$ and let S be a sector. If $\omega$ has $(\omega_2)$ then the class $\widetilde{\mathcal{A}}^{u}_{[\omega]}(S)$ is closed under composition; i.e. for all $f\in\widetilde{\mathcal{A}}^{u}_{[\omega]}(S)$ with $f(0)=0$ and for all $g\in\widetilde{\mathcal{A}}^{u}_{[\omega]}(T)$, where $T\subset \mathcal{R}$ is a sector containing the range of $f$, we have $g\circ f\in\widetilde{\mathcal{A}}^{u}_{[\omega]}(S)$.
	
% \end{theorem}
% \demo{Proof}
% Condition \hyperlink{Mlc}{$(\mathcal{M}_{\on{lc}})$} for $\mathcal{M}_{\omega}$ implies that the map $p\mapsto(W^{(\ell)}_p)^{1/p}$ is non-decreasing for each $\ell>0$, and therefore $\mathcal{M}_{\omega}$ has \hyperlink{pri}{$(\mathcal{M}_{\on{pri}})$}. In addition, if $\omega$ satisfies the condition $(\omega_2)$, then $\mathcal{M}_{\omega}$ satisfies \hyperlink{holom}{$(\mathcal{M}_{\mathcal{H}})$} and therefore \hyperlink{R-Comega}{$(\mathcal{M}_{\{\text{C}^{\omega}\}})$}. Now, it follows from Theorem~\ref{aubsersonmennessiercondition}, arguing for each $\ell>0$, that \hyperlink{Mlc}{$(\mathcal{M}_{\on{lc}})$} implies \hyperlink{salg}{$(\mathcal{M}_{\on{salg}})$} for $\mathcal{M}_{\omega}$. Finally, we can apply Theorem~\ref{Theo.CompStability} and Theorem~\ref{Theo.CompStability.Beur} in order to prove that the class $\widetilde{\mathcal{A}}^{u}_{[\mathcal{M}_{\omega}]}(S)$ is closed under composition, and therefore $\widetilde{\mathcal{A}}^{u}_{[\omega]}(S)$ (thanks to condition $(\omega_1)$).
% \qed\enddemo


\vskip.5cm
\noindent\textbf{Acknowledgements}: The first three authors are partially supported by the Spanish Ministry of Science and Innovation under the project PID2022-139631NB-I00. The research of the fourth author was funded in whole by the Austrian Science Fund (FWF) project 10.55776/PAT9445424.\par


\bibliographystyle{plain}
\begin{thebibliography}{10}
%	\bibitem{BalserUTX} W. Balser, {Formal power series and linear systems of meromorphic ordinary differential equations}, Springer, Berlin, 2000.


%    \bibitem{BrunaInverse} J. Bruna, On inverse closed algebras of infinitely differentiable functions, Studia Math. 69 (1980), no. 1, 59--68.
%
%	\bibitem{FernandezGalbis06}
%	C.~Fernández and A.~Galbis, Superposition in classes of ultradifferentiable functions, Publ. Res. Inst. Math. Sci. 42 (2006), no. 2, 399--419.
%	
%	\bibitem{almostanalytic}
%	S.~F\"{u}rd\"{o}s, D.~N. Nenning, A.~Rainer and G.~Schindl, Almost analytic extensions of ultradifferentiable functions with applications to microlocal analysis,
%	J. Math. Anal. Appl. 481 (2020), no. 1, 123451.
	
	\bibitem{AubersonMennessier81}
	G.~Auberson and G.~Mennessier, Some properties of {B}orel summable functions, J. Math. Phys. 22 (1981), 2472-2481.

	% \bibitem{BraunMeiseTaylor90}
	% R.~W. Braun, R.~Meise and B.~A. Taylor, Ultradifferentiable functions and Fourier analysis, Results Math. 17 (1990), no. 3-4, 206--237.

	\bibitem{ultraholomstab}
	J.~Jim{\'e}nez-Garrido, I.~Miguel-Cantero, J.~Sanz and G.~Schindl, Stability properties of ultraholomorphic classes of {R}oumieu-type defined by weight matrices, Rev. Real Acad. Cienc. Exactas Fis. Nat. Ser. A-Mat. RACSAM 118, 85 (2024).
	
	% \bibitem{index}
	% J.~Jim{\'e}nez-Garrido, J.~Sanz and G.~Schindl, Indices of O-regular variation for weight functions and weight sequences, Rev. R. Acad. Cienc. Exactas Fís. Nat. Ser. A Mat. RACSAM 113 (2019), no. 4, 3659--3697.
	
	\bibitem{sectorialextensions}
	J.~Jim{\'e}nez-Garrido, J.~Sanz and G.~Schindl, Sectorial extensions, via Laplace transforms, in ultraholomorphic
	classes defined by weight functions, Results Math. 74 (2019), no. 1, 27.
	
	\bibitem{sectorialextensions1}
	J.~Jim{\'e}nez-Garrido, J.~Sanz and G.~Schindl, Sectorial extensions for ultraholomorphic classes defined by weight	 functions, Math. Nachr. 293 (2020), no. 11, 2140--2174.
	
%	\bibitem{mixedgrowthindex}
%	J.~Jim{\'e}nez-Garrido, J.~Sanz, and G.~Schindl.
%	block Equality of ultradifferentiable classes by means of indices of mixed
%	{O}-regular variation.
%	block {\em Res. Math.}, 77:art. no. 28, 2022.
	
	\bibitem{Komatsu73}
	H.~Komatsu, Ultradistributions, I. Structure theorems and a characterization, J. Fac. Sci. Univ. Tokyo Sect. IA Math. 20 (1973), 25--105.

%    \bibitem{Krantz02}
%	S. G. Krantz and H. R. Parks, A primer of real analytic functions, 2nd ed., Birkhäuser, Boston, 2002.
%
%	\bibitem{Malliavin} P. Malliavin, Calcul symbolique et sous alg\`ebres de $L^1(G)$, I, Bull. Soc. Math. France 87 (1959), 181--186.

 %    \bibitem{mandelbrojtbook}
	% S.~Mandelbrojt, Séries adhérentes, régularisation des suites, applications, Gauthier-Villars, Paris, 1952.
	
	% \bibitem{matsumoto}
	% W. Matsumoto. Characterization of the separativity of ultradifferentiable classes. J. Math. Kyoto Univ., 24(4):667-678, 1984.

	\bibitem{MozoPhd}
	J.~Mozo-Fernández, Teoremas de división y de Malgrange-Sibuya para funciones con desarrollo asintótico fuerte en varias variables, Ph.D. thesis, University of Valladolid, 1996.

    \bibitem{Poschel}
    J. Pöschel, On the Siegel-Sternberg linearization theorem. J. Dyn. Diff. Equat. 33, 1399--1425 (2021). %https://doi.org/10.1007/s10884-021-09947-7
%	\bibitem{PetzscheVogt}
%	H.-J. Petzsche and D.~Vogt.
%	block Almost analytic extension of ultradifferentiable functions and the
%	boundary values of holomorphic functions.
%	block {\em Math. Ann.}, 267:17--35, 1984.
	
	\bibitem{compositionpaper}
	A.~Rainer and G.~Schindl, Composition in ultradifferentiable classes, Studia Math. 224 (2014), no. 2, 97--131.
	
	\bibitem{characterizationstabilitypaper}
	A.~Rainer and G.~Schindl, Equivalence of stability properties for ultradifferentiable function classes, Rev. R. Acad. Cienc. Exactas Fís. Nat. Ser. A Mat. RACSAM 110 (2016), no. 1, 17--32.
%	
%%	\bibitem{whitneyextensionmixedweightfunction}
%%	A.~Rainer and G.~Schindl.
%%	block On the extension of {W}hitney ultrajets.
%%	block {\em Studia Math.}, 245(3):255--287, 2019.
%	
%	\bibitem{whitneyextensionmixedweightfunctionII}
%	A.~Rainer and G.~Schindl, On the extension of Whitney ultrajets, II, Studia Math. 250(2020), no. 3, 283--295.
%	
	\bibitem{salinas62}
	B.~Rodríguez-Salinas, Clases de funciones analíticas, clases semianalíticas y cuasianalíticas, Rev. R. Acad. Cienc. Exactas Fís. Quím. Nat. Zaragoza (2) 17 (1962), 5--75.
%	
%	\bibitem{RudinDivision} W. Rudin, Division in algebras of infinitely differentiable functions, J. Math. Mech. 11 (1962), no. 5, 797--809.
%
%	\bibitem{Sanzflatultraholomorphic}
%	J.~Sanz, Flat functions in Carleman ultraholomorphic classes via proximate orders, J. Math. Anal. Appl. 415 (2014), no. 2, 623--643.
%	
	% \bibitem{diploma}
	% G.~Schindl, Spaces of smooth functions of Denjoy-Carleman type, Diploma Thesis, Universität Wien, 2009, available online at
	%  \url{http://othes.univie.ac.at/7715/1/2009-11-18_0304518.pdf}.
	
	% \bibitem{dissertation}
	% G.~Schindl, Exponential laws for classes of Denjoy-Carleman differentiable mappings, Ph.D. thesis, Universität Wien, 2014, available online at
	%  \url{http://othes.univie.ac.at/32755/1/2014-01-26_0304518.pdf}.
	
	% \bibitem{testfunctioncharacterization}
	% G.~Schindl, Characterization of ultradifferentiable test functions defined by weight matrices in terms of their Fourier transform, Note Mat. 36 (2016), no. 2, 1--35.
	
%	\bibitem{subaddlike}
%	G.~Schindl, On subadditivity-like conditions for associated weight functions, Bull. Belg. Math. Soc. Simon Stevin 28 (2022), no. 3, 399--427.
%	
%	\bibitem{regularization}
%	G.~Schindl, On the regularization of sequences and associated weight functions, available online at 	 \url{https://arxiv.org/abs/2307.07902}.
%	
%	\bibitem{Siddiqi}
%	J.~A. Siddiqi, Inverse-closed Carleman algebras of infinitely differentiable functions, Proc. Amer. Math. Soc. 109 (1990), no. 2, 357--367.
%	
%    \bibitem{IderSiddiqi}
%	J.~A. Siddiqi and M.~Ider, A symbolic calculus for analytic Carleman classes, Proc. Amer. Math. Soc. 99 (1987), no. 2, 347--350.
%	
%	\bibitem{Thilliezdivision}
%	V.~Thilliez, Division by flat ultradifferentiable functions and sectorial extensions, Results Math. 44 (2003), no. 1-2, 169--188.
	
	\bibitem{thilliez}
	V.~Thilliez, On quasianalytic local rings, Expo. Math. 26 (2008), no. 1, 1--23.
\end{thebibliography}





%\vskip.4cm
%\noindent\textbf{Statements and declarations}:\\
%\vskip.1cm\noindent \textbf{Funding} \ The first three authors are partially supported by the Spanish Ministry of Science and Innovation under the project PID2019-105621GB-I00. The fourth author is supported by FWF-Project P33417-N.\par
%
%\vskip.2cm\noindent \textbf{Competing interests} \ The authors have no relevant financial or non-financial interests to disclose.\par
%
%\vskip.2cm\noindent \textbf{Author contributions} \ All authors contributed to the study and the preparation of previous versions of the manuscript. All authors read and approved the final manuscript.\par
%
%\vskip.2cm\noindent \textbf{Data availability} \ Data sharing not applicable to this article as no datasets were generated or analysed during the current study.

\vskip.5cm
\noindent\textbf{Affiliations}:\\
\noindent Javier~Jim\'{e}nez-Garrido:\\
Departamento de Matem\'aticas, Estad{\'\i}stica y Computaci\'on\\
Universidad de Cantabria\\
Avda. de los Castros, s/n, 39005 Santander, Spain\\
Instituto de Investigaci\'on en Matem\'aticas IMUVA, Universidad de Va\-lla\-do\-lid\\
ORCID: 0000-0003-3579-486X\\
E-mail: jesusjavier.jimenez@unican.es\\

\vskip.1cm
\noindent
Ignacio Miguel-Cantero:\\
Departamento de \'Algebra, An\'alisis Matem\'atico, Geometr{\'\i}a y Topolog{\'\i}a\\
Universidad de Va\-lla\-do\-lid\\
Facultad de Ciencias, Paseo de Bel\'en 7, 47011 Valladolid, Spain.\\
Instituto de Investigaci\'on en Matem\'aticas IMUVA\\
ORCID: 0000-0001-5270-0971\\
E-mail: ignacio.miguel@uva.es\\

\vskip.1cm
\noindent Javier~Sanz:\\
Departamento de \'Algebra, An\'alisis Matem\'atico, Geometr{\'\i}a y Topolog{\'\i}a\\
Universidad de Va\-lla\-do\-lid\\
Facultad de Ciencias, Paseo de Bel\'en 7, 47011 Valladolid, Spain.\\
Instituto de Investigaci\'on en Matem\'aticas IMUVA\\
ORCID: 0000-0001-7338-4971\\
E-mail: javier.sanz.gil@uva.es\\

\vskip.1cm
\noindent Gerhard~Schindl:\\
Fakult\"at f\"ur Mathematik, Universit\"at Wien,
Oskar-Morgenstern-Platz~1, A-1090 Wien, Austria.\\
ORCID: 0000-0003-2192-9110\\
E-mail: gerhard.schindl@univie.ac.at
\end{document}


